% SOURCE docs/release/2026-10-29/source/docs/writing-repair/2026-10-29/src/5b2599cb1b9953cc.tex CHAR 0
\chapter{Assistance from research design to interpretation}
\label{ch:la04-workflow}

The discovery and experiment-design arguments now enter a complete analytical
workflow. This chapter begins with the target and its handoffs, then develops
servicing, branch-rollout and observation examples. Their estimation errors
and policy-version distinctions give concrete tests for the architectural
comparison that follows. The final studies examine data preparation,
validation and reporting, so that a successful operation can be distinguished
from an answer to the bank's original question.

\section{A decision contract before an agent architecture}

A banking assistant becomes useful when it helps an analyst preserve the meaning
of a decision across several transformations. A treasury question becomes a
population and an intervention; account records become measurements of that
population; a diagram becomes an identification problem; an identified
functional becomes a numerical request; and the result becomes a monetary
comparison. Each transformation can change the question without producing a
software error. Fluent explanations are particularly dangerous when they make
these changes hard to notice. The appropriate object of automation is therefore
a sequence of accountable transformations, with explicit unresolved states.

Consider a request to ``reduce the funding cost of the customers likely to
leave.'' An analyst must establish whether the intervention is a deposit-rate
offer, a service action, or a portfolio allocation; whether the unit is an
account, customer, household, or business group; whether likely departure is a
baseline score or a future outcome; and whether cost includes the offer paid to
customers who would have stayed anyway. An assistant can list these alternatives
and attach relevant contract passages. It cannot establish which alternative
the bank intends merely by selecting the most plausible sentence. A structured
interview, conducted without an LLM, is the comparison baseline. The relevant
quality measure is the accuracy and completeness of the resulting target under
expert review, together with analyst time and unresolved-question rates.

We represent the decision contract schematically as
\[
 Q=(\mathcal P,U,A,g_0,g_1,Y,\tau,\nu,\mathcal H,\mathcal A),
\]
where $\mathcal P$ is the population, $U$ the unit, $A$ the feasible action,
$g_0,g_1$ the comparison regimes, $Y$ the outcome, $\tau$ the horizon, $\nu$ the
target measure, $\mathcal H$ the available histories and $\mathcal A$ the declared
assumptions. An ordinary population mean and a propensity-overlap mean have
different $\nu$. A policy responding at day 11 and one acting at day 10 may have
different $\mathcal H$, even when they reference the same event. This tuple is a
project specification, not a theorem that every target is identified. Its
purpose is to make semantic changes observable.

The literature supplies several useful components but no generally qualified
substitute for this contract. Causal discovery assistants can separate user
decisions from formal algorithm outputs; retrieval systems can attach passages;
benchmark environments can require structured identification and estimation
answers. These roles have different failure modes and should not receive a
single undifferentiated ``agent accuracy'' score
\cite{v12_la_2606_23608,v12_la_2402_15301,v12_la_package_causalds}.

\section{Evidence retrieval and the construction of observations}

Retrieval has two distinct products: a discovered source and a supported claim.
A search result supplies the first. The second requires reading primary text,
checking its version and population, and showing why its assumptions apply to
the present query. A banking process manual may establish when a contact occurs
without establishing why customers receive it. A statistical paper may justify
an estimator conditional on exchangeability without establishing that
exchangeability holds in the bank. A source-to-claim relation must therefore
record both the supported statement and the applicability judgment.

The retrieval baseline is a curated collection reviewed by an analyst. A useful
assistant can broaden candidate discovery, locate conflicting passages and
produce an evidence table for review. Evaluation should include relevant,
irrelevant and misleading documents, missing sources, obsolete versions and
sources that contradict each other. Retrieval failure must remain distinguishable
from evidence that a proposed relation is false. The LACR results discussed in
Chapter~\ref{ch:la03-discovery} are a concrete reason not to assume that more retrieved
text monotonically improves discovery \cite{v12_la_2402_15301}.

Observation construction is equally substantive. A field called
\texttt{owner\_id} may identify an operational account holder rather than the
economic owner relevant to exposure aggregation. A link score is a measurement
about a proposed relation, not the relation itself. The observation journey must
retain candidate mappings, inclusion rules, missed records, joint dependence
and the origin of adjudicated labels. Ontology tools can help align terminology,
but class membership and linguistic relations do not identify causal arrows.
The inspected OntoLearner download path also illustrates a reproducibility
problem: the code revision alone does not pin ontology bytes when its hub download
does not specify a revision \cite{v12_la_package_ontolearner}. The design uses
explicit local ontology identities and keeps conceptual variables separate from
their proxies. A real ownership audit remains additional information.

For time-dependent decisions, every relevant measurement has an event time
$t_e$ and an availability time $t_a$. An event observed on day 10 but delivered
on day 11 cannot be used by a day-10 operational policy. Retrospective analysis
may nevertheless use it under a correctly stated design; that is a different
claim from saying a deployed policy could have acted on it. Corrections and
revisions need their own availability records. A validator checks the declared
ordering. Only a process audit can establish whether those declarations reflect
the actual information system.

\section{Diagrams, discovery and identification as separate handoffs}

Diagram elicitation should preserve competing scientific accounts. A proposal
must distinguish a direct edge from an ancestral path, a latent common cause
from a measured covariate, and a circle endpoint from an arrowhead. Revision
history should show which source, reviewer or new experiment changed a claim.
Automatic cycle removal is at most a representation repair: deleting a
scientifically important edge to obtain a DAG does not justify the repaired
mechanism. The ordinary baseline is expert elicitation with a deterministic
graph editor and explicit uncertainty.

Discovery then asks a statistical question under a specified sampling law.
An assistant may suggest a method, summarize its assumptions, or propose a soft
constraint. The numerical procedure remains responsible for its own output
semantics. A response histogram from a language model is not a conditional
independence test under a bank sampling law. A temporal initialization does not
establish that observed anomalies were randomized interventions. A retained
equivalence class should reach the effect-query stage intact. Chapter~\ref{ch:la03-discovery}
explains why language-assisted graph methods require separate tests of
proposal quality, graph retention and downstream effect uncertainty.

Identification maps a target and assumptions to a functional or set. It can
also establish nonidentification in a declared model class. A missing backend,
unsupported graph type or failed calculation instead means that the
implementation did not answer the question. These outcomes are not
interchangeable. The inspected CausalDS helper supports an unknown state and excludes an IV
estimand from nonparametric population-ATE identification. Its result still
requires the criterion scope and error metadata: a negative result can survive
an unavailable fallback, and invalid target nodes can return false with an
error. Section~\ref{exec14:llm-identification} develops that distinction. Its
evaluator also needs care with missing metadata and invalid intervals
\cite{v12_la_package_causalds}. A prototype should preserve that distinction even
when a downstream report would be easier to write as a single yes/no answer.

PyWhy-LLM is a candidate elicitation component. The inspected identification
suggester leaves its front-door routine unfinished and delegates back-door
suggestions to a language model; the separate simple suggester can return
front-door names, as Section~\ref{exec14:llm-identification} explains
\cite{v12_la_package_pywhy_llm}. Such output can propose variables for review; it
does not supply a verified adjustment proof. Similarly, the inspected CLadder
IV template requires generator-level scrutiny before its population-ATE labels
can become benchmark truth \cite{v12_la_package_cladder}. The project example
with half compliers and half never-takers has Wald ratio one and population ATE
two under randomization, exclusion and monotonicity. Correct algebra for the
ratio does not fix an incorrectly named target.

\section{Estimation, diagnostics and the temptation to change the target}

Once identification has supplied a functional, the assistant can prepare a fixed
estimator request, summarize diagnostics and explain the output. It should not
choose a more convenient target after inspecting poor overlap. If trimming
changes the population, that change belongs in a new target or an explicitly
reported sensitivity analysis. An estimated propensity score, nuisance learner,
sample split and covariance rule must be versioned together with the result.
Prediction accuracy does not by itself qualify the resulting effect interval,
especially with dependent observations or adaptively selected specifications.

The inspected OCI example makes the target issue concrete. It averages an AIPW
quantity uniformly over a trimmed sample while using an ATO label
\cite{v12_la_package_oci_agent}. Under the corresponding assumptions this
construction targets the retained population's average effect, not the usual
overlap-weighted population. In the project's two equally frequent strata,
propensities are $0.2$ and $0.5$ and effects are $0$ and $10$. Retaining both
strata gives mean effect $5$. Weighting by $e(X)\{1-e(X)\}$ gives
\[
 \frac{0.16\cdot0+0.25\cdot10}{0.16+0.25}
 =\frac{250}{41}\simeq6.09756.
\]
These are exact arithmetic consequences of different target measures, not a
performance comparison of fitted estimators. They give a small independent
oracle for a semantic validator.

Diagnostics need similarly explicit interpretations. A negative control can
falsify an implication of a causal account; passing it does not prove all
unmeasured confounding absent. A critic can identify a missing source or an
inconsistent formula; agreement between two models exposed to the same text is
not independent scientific confirmation. Uncertainty reporting should preserve
identification uncertainty, sampling uncertainty, sensitivity to assumptions
and numerical approximation separately. A disconnected effect set cannot be
replaced by its midpoint, and a whole-line set cannot be described as a precise
zero effect. Reporting all requested cases, including failures and refusals,
prevents success-conditioned accuracy from masquerading as general availability.

\section{Longitudinal policies require a causal history}

Servicing decisions expose the limits of a generic planner particularly clearly.
Let $L_0$ be baseline information, $A_0$ an initial contact, $L_1$ the subsequent
account state, $A_1$ a second action and $Y$ a later net value. The state $L_1$ can
be affected by $A_0$ while confounding the relation between $A_1$ and $Y$.
Ordinary regression adjustment for $L_1$ can block part of the earlier treatment
effect or induce selection bias. Omitting it can confound the later action.
Hernán and Robins develop the corresponding g-methods and the assumptions under
which they identify longitudinal intervention means \cite{hernan2020}.

For a discrete two-decision example, write
$g_0(l_0)$ for the first action and
$g_1(l_0,a_0,l_1)$ for the second. Under consistency, appropriate sequential
exchangeability and positivity for the policy's supported histories, a
specialization of the g-formula is
\begin{align}
 V(g)=\sum_{l_0,l_1}&
 m\{l_0,g_0(l_0),l_1,g_1(l_0,g_0(l_0),l_1)\}\nonumber\\
 &\quad\times p\{l_1\mid l_0,A_0=g_0(l_0)\}p(l_0),
 \label{eq:la04-policy}
\end{align}
where $m(l_0,a_0,l_1,a_1)=E[Y\mid L_0=l_0,A_0=a_0,L_1=l_1,A_1=a_1]$.
The sum is over supported covariate values; an integral is required for
continuous components. The first treatment changes the distribution over which
the later outcome regression is averaged. That dependence is why simply
conditioning on a fixed observed distribution of $L_1$ can answer a different
question. A recursive regression representation integrates the final outcome
backward, while IP weighting reweights observed histories toward the regime.
Their identification relationship does not imply equal finite-sample behavior.

An assistant can help check whether the proposed regime uses available
information and whether the analyst has stated assumptions at each decision.
It can propose an evidence request for treatment assignment, censoring or action
feasibility. It cannot infer sequential exchangeability from the fact that all
database columns have timestamps. Causally relevant history is not simply every
past variable; some past variables can induce selection bias, as the source's
discussion of history emphasizes \cite{hernan2020}. A policy might also lack
support in distressed segments where staff always take the same action.
More generated explanations do not create that support.

The availability rule is operationally separate from the identification rule.
If the second decision occurs on day 11, a state available on day 11 may enter
$g_1$ under a specified within-day ordering. If the action occurs on day 10, it
may not. A missing within-day ordering must remain unresolved. Censoring and
interference require additional modeling; neither appears automatically in
equation~\eqref{eq:la04-policy}. The assistant's proposed handoff therefore
includes the policy, available history, assignment/retention assumptions,
support diagnostics and explicit limits. This is a design for review, not a
claim that the current orchestration frameworks implement qualified
longitudinal causal inference.

\section{What the available packages can contribute}

The pinned package comparison separates reusable ideas from ready adoption.
PyWhy-LLM offers candidate elicitation interfaces. CausalAgent illustrates a
causal MCP wrapper, but its inspected OLC function is called without being
imported or defined in that module; malformed JSON may become an empty
dictionary in its helper \cite{v12_la_package_causalagent}. That path requires
repair and independent qualification before direct use. The response-based
test in causal\_GPT is not accepted as a bank CI engine, for the reasons derived
in Chapter~\ref{ch:la03-discovery} \cite{v12_la_package_causal_gpt}.

OCI's separation of specifications and reporting is worth evaluating, but its
unbounded notebook-execution route is inappropriate for the proposed fixed-tool
surface. CausalDS is a stronger candidate for a bounded benchmark environment
after scorer and metadata handling are repaired. CLadder remains a useful
formal-question lead with problematic IV cells excluded until reviewed.
OntoLearner may assist semantic mapping with explicit local bytes, but its
extraction assessment also needs the document-identity checks developed in
Chapter~\ref{exec17:ontology}. The adjacent
numerical LCM is neither a language assistant nor a presently qualified temporal
backend: weight identities, component terms and its sufficient lagged model
class need separate treatment \cite{v12_la_package_lcm}.

The Java/Python research bundle LLM-CD and the LLM-SCD notebooks also illustrate
why publication and code adoption need different decisions. Inspected component
licenses do not establish a license for a repository as a whole; missing
root terms and code--paper version alignment remain unresolved
\cite{v12_la_package_llm_cd,v12_la_package_llm_scd}. Their scholarly arguments can
be discussed from genuine papers while code adoption is deferred. No package
in this comparison is credited with a successful local bank or live-model trial.

\section{Orchestration is an implementation choice}

Three generic frameworks were inspected at revisions no later than the client
cutoff. LangGraph supplies state-graph compilation, checkpoint and interrupt
interfaces. Its interrupt documentation explicitly says that resumption starts
the node again, re-executing its preceding logic \cite{v12_la04_langgraph}. A node
that publishes a result before interrupting therefore needs idempotent
publication. A checkpoint is useful state, not proof that an arbitrary external
side effect occurs exactly once. This distinction determines the design of
durable request identifiers and immutable artifact checks.

Pydantic AI supplies typed tool and output boundaries and usage accounting.
In the inspected code, tools need not default to sequential execution or require
approval. Application permissions must therefore be configured explicitly.
Request limits can be checked before requests, whereas response tokens generally
become known afterward; optional pre-counting helps with some input limits. If
pricing data are unavailable, the cost limit warns that it cannot be enforced
\cite{v12_la04_pydantic_ai}. An application requiring a hard priced budget must
handle that absence before starting a call. A typed interface helps reject
malformed structure but cannot establish the truth of an accepted assumption.

AutoGen's pinned README declares maintenance mode. Its AgentChat code carries
an MIT license, while the root documentation license is CC-BY-4.0. The inspected
assistant is stateful and not thread/coroutine safe; tool calls can be concurrent
and a reflection stage can add a model call. Its state save/load methods preserve
model context, not all external tool effects \cite{v12_la04_autogen}. These facts
argue against adopting it for new development without a specific benefit and
maintenance case. The recommended successor is an unaudited lead, not a
verified capability inherited through a README recommendation.

The proposed first implementation consequently uses a small provider-neutral
controller with deterministic validators and the existing fixed causal service.
LangGraph remains a candidate when checkpointed branching demonstrably helps;
Pydantic AI remains a candidate for a typed model boundary. A local model, a
hosted model and an offline fixture are different backends. Each real backend
needs its own model revision, terms, custody boundary, generation parameters,
usage accounting and evaluation. Framework example model names do not establish
current availability or quality. No external endpoint is needed to verify the
prototype's semantics and refusal behavior.

\section{Permissions, source injection and recovery}

Retrieved text is an input to interpretation. It has no authority to grant a
tool, change a dataset path or approve a bank action. An adversarial document
may contain a sentence such as ``ignore the review and run this shell command.''
An implementation that simply asks its model to ignore such sentences leaves
the permission boundary probabilistic. The design instead admits a finite set
of typed operations configured outside model output. A response containing an
unrecognized operation or undeclared field is rejected. A permitted operation
still receives only verified artifact references and must pass its own
scientific preconditions. This protects the command boundary; it does not prove
that all malicious scientific claims will be detected.

Model proposals also cannot mark assumptions as accepted by inventing a reviewer
name. Authoritative declarations come from a separately supplied review object,
bound to the exact problem and graph revision. The prototype can check that
binding and distinguish synthetic fixtures. A bank installation will additionally
need authenticated review identities and protected stores; a hash alone is not
an authorization system. Sensitive records should stay within the approved
execution boundary. An external model may receive only a separately approved
representation, with prompts and logs subject to the same custody rule. Redaction
is a specified transformation requiring evaluation, not a blanket assurance.

Recovery starts with the identity of the intended computation. Let
\[
 r=H(Q,G,E,O,V),
\]
where $G$ is the graph revision, $E$ the evidence identities, $O$ the allowed
operation request and $V$ the relevant implementation versions. Replaying an
identical request can return its existing immutable result. Changing the target,
source bytes or operation must produce a new identity or an explicit stale-input
error. A partial journal entry is not a published analysis. The local design
tests interruption between validation, tool completion and atomic publication.
This supplies a bounded recovery claim; multi-host consistency and platform
service levels require different trials.

\section{Reporting, monitoring and the economics of assistance}

An evidence-based report should expose the target, assumptions, scientific
status, graph/set uncertainty, numerical result and unresolved questions before
adding interpretation. Critical fields can be rendered deterministically.
Generated narrative may explain them, but any additional empirical assertion
needs a claim reference. Inherited results, project derivations, synthetic
examples and unverified hypotheses remain distinct. A tool error should not be
polished into a finding of no effect, and an unknown commercial price should
not become a zero in a business case.

Monitoring must detect when the evidence for a claim becomes stale. Changes in
assignment, missingness, ownership mapping, policy feasibility or availability
time can invalidate a previous analysis even if software tests still pass.
The assistant may triage an alert and propose an investigation. Suspension and
resumption rules belong to the analysis owner. Retuning prompts or refitting a
learner after inspecting drift creates a new procedure whose old calibration
does not transfer automatically.

The economic comparison includes model calls, retrieval, numerical computation,
review, failures, rework and ongoing maintenance. If $C_m,C_r,C_n,C_h,C_f$ denote
these components, a proposed net benefit is a change in held-out decision value
minus $C_m+C_r+C_n+C_h+C_f$, relative to a specified ordinary workflow.
This expression is an accounting design, not an ROI estimate. The available
evidence does not supply authentic bank losses, prices, review times or accepted
decision outcomes. A faster explanation can be valuable, but it must not conceal
a lower rate of scientifically usable answers. The appropriate next trial
therefore measures decision quality, all-request availability and total cost
separately, while the implementation preserves the reasons an answer was
ambiguous, refused or unavailable.

The following developments make those handoffs testable. Servicing, branch
adoption and observation construction require different intervention and
sampling laws. Their exact examples show what a workflow must preserve before
the architectural alternatives and final validation study can be assessed.


% SOURCE docs/release/2026-10-29/source/docs/writing-repair/2026-10-29/src/fbed90a4ad9bc3f8.tex CHAR 0
\section{Servicing decisions: intervention histories and net value}
\label{exec14:l5-servicing}

The policy functional in equation~\eqref{eq:la04-policy} becomes operationally
useful only when its histories have a banking interpretation. Consider an
initial servicing contact $A_0$, a subsequent distress indicator $L_1$, and a
second contact $A_1$. An initial contact can prevent distress; distress can
prompt a second contact; and both contacts incur costs. A request to compare
``contacted and uncontacted customers'' conceals all three relations. It also
conceals whether the analyst wants the value of a fixed history, a rule that
responds to distress, or a policy learned from the same observations. These are
different inferential requests. An assistant can make the alternatives explicit
and translate the selected request into a rule. Its interpretation remains
reviewable against the business decision.

The following mechanism is stipulated for an independent exact reference; its
numbers are illustrative USD per account, not observations from a bank
\cite{exec14_workflow_oracle}. There is no baseline covariate in this small law.
The first action has probability $1/2$, and the state transitions satisfy
\[
 P(L_1=1\mid A_0=0)=\frac45,\qquad
 P(L_1=1\mid A_0=1)=\frac25.
\]
The observational second-action probabilities are
\[
 e_1(l)=P(A_1=1\mid A_0,L_1=l)
 =\begin{cases}1/2,&l=0,\\3/4,&l=1.\end{cases}
\]
Gross recovery is $100-40L_1+10A_0+30L_1A_1$. Each contact costs eight, so the
outcome to be compared is
\[
 Y=100-40L_1+2A_0+(30L_1-8)A_1.
 \label{eq:exec14-l5-net}
\]
The stipulated mechanism has independent assignment noise at each action,
conditional on the displayed history, and a deterministic final outcome.
This supplies the exchangeability needed for its illustrative intervention
interpretation. An observed bank table with the same conditional means would
still require an assignment argument; numerical resemblance would not supply
that premise. Treatment-confounder feedback and the causal conditions of the
g-formula are developed by Hern\'an and Robins \cite{hernan2020}.

Compare three fixed, feasible policies. The no-contact policy sets both
actions to zero. The treat-all policy sets both to one. The risk-responsive
policy sets $A_0=1$ and $A_1=L_1$. Its second contact is affordable and available
only after the distress state is delivered. The policy does not use a future
outcome or a revised state unavailable at that decision. Under the risk policy,
net recovery is 102 when $L_1=0$ and 84 when $L_1=1$. Averaging under its own
post-first-action state law gives
\[
 V(g_{\mathrm{risk}})=\frac35(102)+\frac25(84)
 =\frac{474}{5}=94.8.
\]
The corresponding values are $V(g_{\mathrm{none}})=68$ and
$V(g_{\mathrm{all}})=90$. The risk-responsive policy gains
$134/5=26.8$ over no contact and 4.8 over contacting everyone. The first contrast
decomposes into two from the first contact's direct net contribution, sixteen
from the reduction in distress probability, and $44/5=8.8$ from the second
contact's net contribution under the changed state distribution. This
decomposition explains why a recovery lift alone is insufficient for the
business decision. A contact can improve gross recovery while its marginal
net contribution is negative in a low-risk state.

Two plausible shortcuts answer different questions. Restricting the observed
data to histories that happen to follow the risk policy gives mean 93, because
the two retained states have equal conditional frequency. Restricting to
histories with both actions zero gives $220/3$, because distress changes the
chance of avoiding the second contact. The unweighted difference is therefore
$59/3$, rather than $134/5$. Alternatively, averaging the risk-policy outcome
under the pooled distress frequency $3/5$ yields 91.2. That calculation holds
the mediator distribution at its observational mixture, rather than allowing
the first intervention to change it. Neither error is repaired by producing
an explanation of a correctly computed mean. The assistant's most valuable
intervention is to detect the target substitution before numerical execution.

For a general baseline $L_0$, write
$H_1=(L_0,A_0,L_1)$ and
$m(H_1,a_1)=E[Y\mid H_1,A_1=a_1]$. With a fixed regime $g$, define
\begin{align*}
 m_g(H_1)&=m\{H_1,g_1(H_1)\},\\
 Q_0(L_0)&=E[m_g(H_1)\mid L_0,A_0=g_0(L_0)].
\end{align*}
Then $V(g)=E[Q_0(L_0)]$ under the stated longitudinal identification
conditions. Backward outcome regression implements these successive
conditional expectations. Forward g-computation models the state transition
and outcome, then propagates the regime through both. Inverse-probability
weighting instead uses
\[
 w_0=\frac{\mathbf1\{A_0=g_0(L_0)\}}
 {P\{A_0=g_0(L_0)\mid L_0\}},\qquad
 w_1=\frac{\mathbf1\{A_1=g_1(H_1)\}}
 {P\{A_1=g_1(H_1)\mid H_1\}},
\]
for which $E[w_0w_1Y]=V(g)$. All three representations use the intervention's
history. They differ in which nuisance quantities must be estimated and in
their behavior when data are sparse, weights are variable, or a regression
extrapolates. The choice should follow the declared target and its supported
histories, rather than the name of a preferred package.

\section{Retention, recursive augmentation and uncertainty}
\label{exec14:l5-retention}

The outcome may become unavailable when an account moves to another servicing
system or its case is unresolved. Let $R$ indicate that $Y$ is observed at the
final horizon. To identify the same full-population value, suppose
$R\perp Y\mid H_1,A_1$, with
$r(H_1,A_1)=P(R=1\mid H_1,A_1)>0$ on the regime's support. This is a final-outcome
missingness law. If earlier histories also disappear, their observation
mechanisms require additional nodes and weights; the single $R$ does not
silently represent them. Nor does administrative censoring justify this
conditional independence merely because its date is recorded
\cite{hernan2020}.

In the stipulated servicing law, let retention be one for $L_1=0$ and one half
for $L_1=1$, independently of the outcome given history. Among retained accounts
under the risk intervention, the mean becomes
\[
 \frac{(3/5)102+(2/5)(1/2)84}{3/5+(2/5)(1/2)}=97.5.
\]
Inverse retention weights restore the intervention value 94.8. Multiplying
$w_0w_1$ by $R/r$ also restores it in the observational enumeration. The
complete-case mean is higher because distressed accounts are less often
retained, not because servicing became more effective. A fluent report that
describes 97.5 as a population policy value would obscure an observation-law
error. Conversely, if loss depends on an unobserved recovery component even
after conditioning, these weights have no such causal justification. The
appropriate output is a sensitivity analysis or unresolved population value,
not a fabricated retention probability.

An augmented representation shows where the pieces enter. With the true
nuisance quantities, a centered score for this two-action, final-retention
problem is
\begin{align}
 \phi_g(O)={}&w_0w_1\frac{R}{r(H_1,A_1)}
          \{Y-m(H_1,A_1)\}\nonumber\\
 &+w_0\{m_g(H_1)-Q_0(L_0)\}+Q_0(L_0)-V(g).
 \label{eq:exec14-l5-score}
\end{align}
The first residual has conditional mean zero under the observation law. The
second has mean zero conditional on $L_0$ after first-action reweighting.
The last term centers the baseline regression. Where $R=0$, the first term
is defined as zero; it does not require an unavailable value of $Y$. This
mean-zero argument is a population identity. A fitted estimator needs control
of nuisance errors and remainder terms, a valid sampling law, and the
appropriate handling of sample splitting. The display does not establish
coverage simply by being called an influence function.

The exact law isolates one aspect of precision. Its outcome is deterministic
given history and it has no varying $L_0$. The first and last components in
equation~\eqref{eq:exec14-l5-score} vanish, leaving variation from the intermediate
state. Under independent accounts, the risk-policy score has variance
$3888/25=155.52$; the no-contact score has variance 512. Their covariance is
zero here because their nonzero first-action indicators are disjoint and both
scores are centered. The variance of their difference is consequently
$16688/25=667.52$. A normalized policy-IPW estimator with true weights has
asymptotic centered score $w_0w_1\{Y-V(g)\}$ when there is no retention loss.
Its variance is 248.832 for the risk policy and 1228.8 for no contact.
These are exact moments under the declared law
\cite{exec14_workflow_oracle}. They explain how recursive adjustment can remove
second-action sampling variation in this example. They do not rank fitted
procedures under arbitrary bank dependence or qualify any confidence interval.

Robustness requires equally precise terminology. Luedtke and colleagues
distinguish conventional longitudinal double robustness, requiring all outcome
components or all treatment components to be correct in the relevant sense,
from sequential double robustness, allowing the appropriate component to be
correct at each stage \cite{sequentialdr2018}. Their definitions distinguish
outcome-regression specification from treatment-mechanism consistency. A
point-treatment AIPW function placed inside a loop does not inherit their
recursive construction or its theorem. An assistant recommending a
``doubly robust longitudinal estimator'' should therefore identify the actual
algorithm and its conditions, rather than treating robustness as a universal
package property. The safe comparison includes a transparent fixed-regime
g-formula, an IP-weighted representation, and an estimator whose longitudinal
construction is genuinely supported by primary text.

Support is a property of the policy's histories, not only a histogram of the
first propensity. An initially supported contact can lead to a distressed
state where the requested later action is never taken. Product probabilities
can also become very small although each decision separately looks tolerable.
Weight truncation may improve numerical behavior while changing bias and the
effective comparison. Diagnostics should show where weights concentrate,
which histories disappear, and how the result changes under a declared
sensitivity choice. A regime that requires a forbidden or unavailable contact
should be returned for redesign. A new policy learned after inspecting the
same outcomes introduces selection and policy-learning questions beyond the
fixed-policy arithmetic above. Its apparent best value cannot be reported
with a fixed-regime interval without a separate argument.

\section{Branch adoption is a panel question with a different target}
\label{exec14:l5-panel}

Suppose branches adopt a servicing platform at different dates and the bank
asks whether the platform improves recovery. This differs from evaluating the
account-level dynamic contact rule. Let $G_i$ be branch $i$'s first adoption
time, $Y_{it}(g)$ its outcome at time $t$ if adoption begins at $g$, and
$Y_{it}(\infty)$ the outcome without adoption. A useful target is
\[
 \operatorname{ATT}(g,t)
 =E[Y_{it}(g)-Y_{it}(\infty)\mid G_i=g],\qquad t\geq g.
\]
The unit, adoption event, composition of branch accounts, and monetary outcome
must be fixed before estimation. A decline in average loss among remaining
accounts could arise from account transfers rather than an improvement for
the original cohort. A repeated cross-section and a stable account panel
therefore need different observation arguments. Arkhangelsky and Imbens
organize panel methods by data configuration, relative numbers of units and
periods, and treatment assignment patterns; this organization helps avoid
selecting a method solely because the table has a date column
\cite{panelreview2024}.

With no anticipation and an appropriate never-treated comparison group $C$,
unconditional parallel untreated trends over the relevant periods identify
\begin{align*}
 \operatorname{ATT}(g,t)
 ={}&E[Y_{it}-Y_{i,g-1}\mid G_i=g]\\
 &-E[Y_{it}-Y_{i,g-1}\mid C_i=1].
\end{align*}
Conditional parallel trends instead require comparing changes at the relevant
baseline covariates $X$ and averaging over the adopting cohort's covariate
distribution, with overlap for that comparison. Callaway and Sant'Anna also
develop not-yet-treated comparisons and explicit anticipation windows
\cite{callaway2021}. For an anticipation horizon $\delta$, the baseline must
precede $g-\delta$, and a comparison group must still be unaffected at the
relevant later time including its anticipation window. A branch preparing for
rollout can change staffing before formal adoption; the recorded software
installation date alone does not establish no anticipation.

The comparison choice carries substantive assumptions. Never-treated branches
may operate in different markets. Not-yet-treated branches may be more similar
but supply no untreated comparison after the last adoption, and their rollout
plans may already affect behavior. Both designs require an argument about
untreated changes, rather than the stronger sequential exchangeability of an
account-level dynamic policy. The analyst must state which argument is being
used. An assistant can locate adoption notices, reconstruct candidate dates,
and propose competing controls. It cannot establish the unobserved untreated
trend from a satisfactory regression fit.

A small stipulated panel makes these distinctions visible
\cite{exec14_workflow_oracle}. There are three equal-sized cohorts, early
adopters $E$, late adopters $L$, and never adopters $N$. Their observed means are
\begin{center}
\begin{tabular}{lrrr}
\toprule
Cohort & Time 0 & Time 1 & Time 2\\
\midrule
Early &100&107&119\\
Late &100&102&108\\
Never &100&102&104\\
\bottomrule
\end{tabular}
\end{center}
The stipulated untreated means are 100, 102, and 104 for every cohort.
Early adoption has effects five at time 1 and fifteen at time 2; late adoption
has effect four at time 2. Comparing early and never cohorts from time 0 to
time 1 gives five; from time 0 to time 2 gives fifteen. Comparing late and
never cohorts from time 1 to time 2 gives four. In contrast, comparing the
late cohort's change to the already-treated early cohort's change gives
\[
 (108-102)-(119-107)=-6.
\]
The early cohort's treatment effect increases by ten during that comparison.
The negative difference therefore subtracts another group's changing effect
from the late cohort's benefit. It is not evidence that the platform harms
late adopters.

Aggregation also changes the question. Equal weighting of the three treated
cohort-time cells gives $(5+15+4)/3=8$. Equal weighting of the two cohorts at
their first exposure gives $(5+4)/2=4.5$. The first includes a later-exposure
effect for early adopters; the second compares a common event horizon. Other
bank questions may weight cohort sizes, account exposure, or the distribution
of dates in a proposed rollout. These weights must be stated. In this particular
balanced fixture, residualizing the adoption indicator against cohort and
calendar effects produces a two-way fixed-effects coefficient of 4.5: the
early time-2 cell has zero residual weight and the other treated cells have
equal positive weights. It happens to equal the first-exposure target here.
The contaminated two-group comparison does not imply that every fixed-effects
coefficient is negative or scientifically useless. The exact residualization
and the survey's qualified discussion both argue for inspecting the target and
weights, rather than adopting a universal rejection of fixed effects
\cite{panelreview2024,exec14_workflow_oracle}.

Observational ambiguity remains even after this algebra is correct. Suppose
the early cohort's untreated outcome at time 2 is 107 because it receives a
cohort-specific favorable shock of three, and its actual treatment effect
is twelve. The observed value remains 119, and every other displayed value is
unchanged. The first law's effect fifteen and the alternative law's effect
twelve are compatible with the same observed panel. The alternative violates
the relevant post-adoption parallel untreated trend. Additional common
pre-adoption periods can be appended to both laws without distinguishing
them. Good pre-adoption agreement is a useful diagnostic; it does not exclude
every later counterfactual shock. A bank may need a rollout assignment design,
a defensible control, or explicit sensitivity to such shocks before attaching
a causal interpretation to the contrast.

Synthetic control and factor-based alternatives are substantive when the
bank has few treated units, rich histories, and a credible structure for
untreated outcomes. They do not automatically solve an unsupported
counterfactual: donor contamination, a new factor after rollout, or a change
in measurement can defeat their argument. The panel survey motivates these
alternatives without asserting a particular bank guarantee
\cite{panelreview2024}. The assistant's proposal should connect each candidate
to the data dimensions and assignment pattern, identify which restriction
supplies its counterfactual, and retain uncertainty about that restriction.
A dynamic servicing value should not be silently substituted for a branch
adoption ATT merely because both methods return a monetary number.

\section{Observation construction can determine the answer}
\label{exec14:l5-observations}

The preceding values assume that accounts and outcomes have been correctly
assembled. A causal workflow should instead begin with the population frame,
inclusion mechanism, ownership relation, and label meaning. An assistant
extracting ``recovery'' from notes may confuse gross recovery with net value,
or a revised figure with the figure available at the decision. Linking two
accounts to a customer may change the unit, create duplicated outcomes, and
alter dependence at once. These changes belong to the scientific observation
model. A schema that accepts the resulting rows does not establish their truth.

Consider a frame containing two equally frequent strata with values zero and
100. Suppose every zero-value unit is included but each 100-value unit is
included independently with probability $1/4$. Writing $S$ for inclusion,
the mean among observed records is
\[
 \frac{E[SY]}{E[S]}=
 \frac{(1/2)(1/4)100}{(1/2)1+(1/2)(1/4)}=20.
\]
The population mean is 50. If the frame size $N$ and truthful inclusion
probabilities $\pi_i$ are known, the Horvitz--Thompson form
\[
 \widehat\mu_{\mathrm{HT}}=
 \frac1N\sum_{i=1}^N\frac{S_iY_i}{\pi_i}
\]
has expectation 50 under the stipulated sampling law. This uses the known
frame denominator; a ratio estimator with an estimated denominator is a
different construction. For independently generated framed units, the
per-unit weighted contribution is zero or 400. Its variance is
$20000-50^2=17500$, so the mean's variance is $17500/N$
\cite{exec14_workflow_oracle}. Correcting selection can therefore increase
uncertainty even while restoring the target. A narrow unweighted interval
around 20 does not constitute a more precise estimate of 50.

The correction requires knowledge that an LLM cannot manufacture. If inclusion
probabilities are estimated from an incomplete frame, depend on an unknown
outcome component, or are zero for a relevant group, the displayed equality
need not apply. Retrieving a description of the sampling system can identify
a candidate probability model; an actual audit must establish that it was
followed. When the target is an intervention mean, treatment, retention,
sampling, and linkage cannot simply be handled as unrelated cleaning steps.
Their joint chronology determines which histories and outcomes are observed
and which probabilities can be conditioned on. The longitudinal observation
intervention discussed by Hern\'an and Robins supplies a model-based way to
formulate that joint problem \cite{hernan2020}.

Noisy labels pose a separate problem. Let $T$ be a true binary ownership or
distress label, $Z$ its observed classifier label, $p=P(T=1)$, sensitivity
$s=P(Z=1\mid T=1)$ and specificity $c=P(Z=0\mid T=0)$. If these error rates are
known for the target population, then
\[
 P(Z=1)=sp+(1-c)(1-p),\qquad
 p=\frac{P(Z=1)+c-1}{s+c-1}.
\]
With $p=0.4$, $s=0.8$ and $c=0.9$, the positive-label frequency is 0.38 and
the inversion returns 0.4. If $s+c=1$, the observed label frequency contains
no information about prevalence through this model. Close to that boundary,
sampling noise is amplified by the denominator. With known rates and $n$
independent labels in this example, the inverted proportion has variance
$(589/1225)/n$ before any probability-bound handling
\cite{exec14_workflow_oracle}. Estimating the error rates from a gold-label
sample adds uncertainty. Reusing one error-rate estimate across many account
records also introduces shared estimation error; multiplying the number of
rows does not make that component disappear.

This inversion concerns a marginal prevalence, not a causal effect or an
individual ownership map. Errors can differ by treatment, history, market or
time, and two competing customer mappings can have the same correctness
count. A prevalence correction does not reconstruct which record belongs to
which economic owner. For causal use, the workflow must retain the candidate
mapping and how its uncertainty affects exposure, outcome and unit formation.
An LLM can propose label rules and flag contradictory notes; the proposal
should compare that assistance with deterministic rules and human adjudication
on genuinely reviewed examples. Agreement between two assistants exposed to
the same notes is not an independent gold label. Missing truth is an evidence
limitation, rather than a reason to replace adjudication with confidence in
generated prose.

Dependence can survive correct inclusion and label handling. Let $B$ independent
branches each have a centered shock $\varepsilon_b$ of variance one, shared
by two account records. Their overall account mean is
$B^{-1}\sum_b\varepsilon_b$, with variance $1/B$. Treating the $2B$ rows as
independent instead gives $1/(2B)$. Now add a centered calendar shock $C$ of
variance four shared by every branch, independently of their shocks. The
mean's variance is $4+1/B$. More accounts or branches reduce only the second
component. This is an exact covariance construction
\cite{exec14_workflow_oracle}, not a universal banking variance model. It
shows why account counts alone cannot justify precision and why a common
calendar component needs replicated dates or additional restrictions.
Splitting nuisance training and evaluation by independent branch can protect
one dependence boundary; it cannot remove the shared calendar shock. The
sampling and covariance argument must follow the intended deployment and
observation law.

\section{What assistance and software reuse can responsibly achieve}
\label{exec14:l5-handoffs}

These three pathways suggest distinct roles for an assistant. In sequential
servicing, it can translate a verbal rule into a decision-time history, locate
post-action variables, and ask why assignment and retention are conditionally
ignorable. In a branch rollout, it can distinguish cohort time from event time,
construct candidate comparison groups, and expose anticipation or account
composition changes. In observation construction, it can extract candidate
definitions and mappings while retaining conflicting evidence and unavailable
truth. A generic request to ``run causal inference'' is insufficient for all
three. The proposed product should make the pathway and its resulting target
visible to the analyst before selecting an estimator.

The point-in-time OCI workflow demonstrates a narrower literature precedent:
an accountable human principal, actor/critic assistance, reproducible notebooks,
cross-fitted AIPW and overlap diagnostics under supplied unconfoundedness
\cite{exec13_data_oci}. It supports discussion of human oversight and
point-treatment workflow. It does not establish sequential exchangeability,
identify a rollout counterfactual, or qualify ownership reconstruction. These
are candidate extensions requiring the corresponding scientific and numerical
arguments. A reviewer should be able to see whether the assistant asks for
the missing premise, changes the target explicitly, or refuses the requested
interpretation. An apparently complete notebook can still be wrong if its
target differs from the business question.

The inspected R package \texttt{ltmle} offers a relevant longitudinal
component. It estimates means specific to treatment and censoring, supports
dynamic treatment rules, and returns outcome-regression and weighting
outputs \cite{sw_ltmle}. Its interface
uses column order to encode model chronology and binary treatment nodes.
An adapter would need to validate the actual event and availability clocks
before constructing that order. A continuous contact dose cannot be passed
as though it were binary without an explicit intervention redesign.
Absorbing event coding is appropriate for a one-time event only when its
semantics are actually absorbing. Automatically truncating net outcomes at a
declared range can alter the quantity being averaged, whereas a known affine
scaling can be reversed for an intervention mean. These interface choices
have scientific consequences, rather than merely formatting consequences.
The saved source inspection motivates restricted reuse; it is not an execution
or qualification of the package in a bank setting.

A matched baseline is an analyst using explicit decision interviews, a graph
editor, and deterministic estimation software with the same sources and
information budget. Assistance should be assessed on preserved targets,
unsupported assumptions, actual tool completion, accurate uncertainty, and
review effort. Comparisons should include unavailable histories, unsupported
actions, contaminated controls, missing frames, and labels without gold truth.
An appropriate unresolved answer may be more valuable than a confidently
generated estimate. Genuine improvement would mean reducing analyst effort
while preserving those decisions and making more errors detectable, rather
than increasing the number of successful-looking reports.

The immediate prototype can use the finite mechanisms above and suitable
public panels to test those transformations economically. Known synthetic
truth tests target preservation and numerical arithmetic; a public observed
panel tests the discipline of stating a counterfactual without supplying its
truth by construction. A later bank experiment must resolve the actual
assignment, observation, operating-cost and acceptance questions. The proposal
can justify that experiment through its scoped comparisons and explicit
failure modes. It cannot replace the experiment's missing evidence with the
existence of an agent architecture.
\label{exec14:l5-development-end}


% SOURCE docs/release/2026-10-29/source/docs/writing-repair/2026-10-29/src/e567d0206e1efd7f.tex CHAR 0
\section{How estimated components can change a sequential value}
\label{exec14:l5-remainder}

The exact score moments above assume the true conditional quantities. A
prototype must also make visible what happens when its fitted components are
wrong. This is particularly relevant when an assistant chooses a regression,
uses an estimated assignment model, and presents the resulting augmented
calculation as automatically protected against bias. Double robustness is a
property of a specified construction under specified conditions. It is not a
guarantee that any two imperfect components cancel each other.

Consider the same fixed risk policy, without retention loss or baseline
variation. Let $p_0$ be the observational probability of its first action and
$p_1(l)$ the probability of its required second action at state $l$. These are
probabilities of following the rule: when the rule requests zero, $p_1(l)$ is
the probability of action zero. Let $\widetilde p_0$ and
$\widetilde p_1(l)$ be fixed positive approximations, and define
\[
 \rho_0=\frac{p_0}{\widetilde p_0},\qquad
 \rho_1(l)=\frac{p_1(l)}{\widetilde p_1(l)}.
\]
Write $\widetilde m_g(l)=m_g(l)+d(l)$ and
$\widetilde q=V(g)+\eta$. The last quantity is an approximation to the
first-stage regression; it need not equal the average of
$\widetilde m_g$ under the true transition law. Replace the quantities in
equation~\eqref{eq:exec14-l5-score} by these fixed approximations and add
$V(g)$ to obtain an uncentered augmented value $\widetilde Z$. Conditioning
successively on the action history gives
\begin{align*}
 E[\widetilde Z]
 ={}&-\rho_0E_g[\rho_1(L_1)d(L_1)]\\
 &+\rho_0E_g[m_g(L_1)+d(L_1)-\widetilde q]
   +\widetilde q.
\end{align*}
Here $E_g$ averages over the state law after the regime's first action.
Subtracting $V(g)$ yields the exact identity
\begin{equation}
 E[\widetilde Z]-V(g)
 =(1-\rho_0)\eta
   +\rho_0E_g[(1-\rho_1(L_1))d(L_1)].
 \label{eq:exec14-l5-remainder}
\end{equation}
This is a project derivation for this restricted two-stage law, independently
checked by enumerating the observational histories
\cite{exec14_workflow_remainder}. It does not claim the general sequential
robustness theorem discussed by Luedtke and colleagues
\cite{sequentialdr2018}.

The identity separates two errors with different histories. Correct action
probabilities at both stages make the right side zero. Correct outcome
quantities at both stages also make it zero. Other combinations can make it
zero, but only by satisfying the displayed stage conditions. If $d=0$ while
both $\rho_0\neq1$ and $\eta\neq0$, correctness of the final outcome model
alone does not remove the first-stage product. Conversely, a correct
first-action probability does not remove a second-stage treatment/outcome
product. Fitting two separate models and calling one of them a critic has no
bearing on these conditions.

For a concrete failure, use the servicing law's $p_0=1/2$, approximate it by
$\widetilde p_0=2/5$, and set $\widetilde q=V(g)+3$. In the low-distress
state the rule requires action zero with true and approximated probability
$1/2$. In the distressed state it requires action one with true probability
$3/4$ and approximated probability $1/2$. Let $d(0)=0$ and $d(1)=5$.
Then $\rho_0=5/4$, and the second product averaged under the regime is
$(2/5)(1-3/2)5=-1$. The total bias is
\[
 (1-5/4)3+(5/4)(-1)=-2.
\]
The augmented mean becomes 92.8 instead of 94.8. Enumeration of the weighted
residual, intermediate regression and final scalar separately gives the same
answer. Restoring both assignment probabilities restores 94.8 despite these
outcome errors; restoring both outcome components also restores it despite
the assignment errors. These are population identities with fixed
approximations, not fitted-estimator performance results
\cite{exec14_workflow_remainder}.

There is also a quantitative diagnostic. With square-integrable components
under the regime's state law, Cauchy--Schwarz gives
\[
 |E[\widetilde Z]-V(g)|
 \leq |1-\rho_0|\,|\eta|
 +\rho_0\|1-\rho_1\|_{2,g}\|d\|_{2,g}.
\]
Lower bounds on approximated treatment probabilities are relevant because
they control the ratios; weak support can amplify small absolute propensity
errors. The bound does not supply observable values of the unknown errors.
It explains which approximations must be justified, and why predictive fit
alone need not settle the causal quantity.

For nuisance quantities trained on an independent sample, the fixed-component
calculation can be read conditionally on that training sample. A familiar
route to root-$n$ inference requires the bias products to be
$o_p(n^{-1/2})$, appropriate convergence of the centered score, and a valid
limit law for its sampling unit. Each requirement needs its own argument;
sample splitting alone proves none of them. Under the independent-account,
known-law construction above, the contrast's standard-error scale is
$\sqrt{667.52/n}$. A normal planning multiplier of 1.96 would require at
least 4007 accounts to make its planning half-width no greater than USD 0.8.
This is arithmetic using the known-law variance, not a claimed achieved
confidence interval, a selected trial size, or a bank precision commitment
\cite{exec14_workflow_remainder}. Dependence, estimated components and rare
histories can invalidate that planning scale. Increasing the account count
does not remove the fixed bias of two in the preceding example.

The assistant should therefore keep the causal target, fitted components,
training boundary, covariance unit and approximation status together. It can
flag a mistaken probability of action one when the policy requests zero,
propose checks of each stage, and compare a transparent fixed-regime
calculation. An unavailable rate argument remains an uncertainty limitation
even when the numerical value is finite. A report should explain that
limitation before turning the point estimate into a recommendation.

\section{Conditional panel comparisons must preserve cohort composition}
\label{exec14:l5-composition}

The preceding panel example used identical untreated trends across cohorts.
Conditional parallel trends allow a different situation: branches with the
same baseline covariates have comparable untreated changes, while adopting
and control branches have different covariate compositions. Pooling their
changes without standardization can then answer the wrong question even
though the conditional design is valid. Callaway and Sant'Anna's
identification representations average the conditional comparison over the
adopting group's covariate distribution \cite{callaway2021}.

Here is an independent stipulated two-period example in USD per account
\cite{exec14_workflow_remainder}. Let $X$ distinguish two branch markets,
and suppose three quarters of adopting branches have $X=1$, compared with
one quarter of never-adopting controls. At $X=0$, the untreated change is two
and the adoption effect is four. At $X=1$, those quantities are ten and
eight. Within each market the untreated change is identical for adopting
and control branches; no anticipation and stable composition are stipulated.
The target is the adopting cohort's average treatment effect,
\[
 \operatorname{ATT}=\tfrac14(4)+\tfrac34(8)=7.
\]
The adopting cohort's observed average change is
$\tfrac14(2+4)+\tfrac34(10+8)=15$. The pooled control change is
$\tfrac34(2)+\tfrac14(10)=4$. Their unadjusted difference is eleven,
overstating the target by four. The difference is caused by unequal market
composition, not by a violation of the stipulated conditional trends.

Outcome standardization instead computes the control change under the
adopting composition: $\tfrac14(2)+\tfrac34(10)=8$. Subtracting it from
fifteen restores seven. Weighting supplies a second representation. Suppose
the two cohorts have equal population mass, so the propensity to belong to
the adopting cohort is $e(0)=1/4$ and $e(1)=3/4$. Control weights proportional
to $e(X)/\{1-e(X)\}$ are $1/3$ and three. Their control-group average is
one, and their weighted change is
\[
 \tfrac34\left(\tfrac13\right)2
 +\tfrac14(3)10=8.
\]
The weights reconstruct the adopting covariate distribution. They are
cohort-comparison weights, rather than the probabilities of following a
two-action account policy. The corresponding target and observation clocks
must remain distinct when an assistant selects an API.

If market $X=1$ has no untreated controls, this representation cannot be
evaluated from the stipulated comparison without extrapolation or another
restriction. Dropping that market changes the target from seven to four;
silently describing the retained-market result as the original ATT would
replace the business question. Matching, regression and weighting should
therefore be compared on their supported target, not merely on whether each
returns a coefficient. Adding a covariate measured after platform adoption
can also change the scientific argument; the baseline market variable here
does not justify adjustment for any convenient post-adoption metric.

For an illustrative deployment cost of six per adopting account, the
identified gross effect of seven leaves a net margin of one. The pooled
difference of eleven would suggest five. These are planning values under
stipulated truth and cost, with no fitted uncertainty or bank ROI claim
\cite{exec14_workflow_remainder}. They show why composition and support checks
matter to the investment proposal: a target error can dominate a modest
economic margin even when every arithmetic step and generated explanation
appears convincing.
\label{exec14:l5-depth-end}


% SOURCE docs/release/2026-10-29/source/docs/writing-repair/2026-10-29/src/85d18da2dea7b7c7.tex CHAR 0
\section{From language-assisted collaboration to adaptive servicing}
\label{exec16:longitudinal-policy}

A bank considering language-assisted servicing faces two related decisions.
An analyst may ask an assistant to formulate and evaluate a contact policy;
the contact itself may also contain model-generated language. In the first
case, the assistant changes the analysis workflow. In the second, it changes
the intervention delivered to a customer. A successful analysis assistant
does not establish the effectiveness of the resulting message, and an
effective writing strategy does not establish a sound policy evaluator.
The distinction determines what the proposed product should estimate, which
public evidence is relevant, and what an inexpensive prototype must resolve.

The sequential servicing construction in
Section~\ref{exec14:l5-servicing} already gives a formal baseline for a
fixed policy with treatment--confounder feedback.
Here the question is whether the available language-model literature and
software preserve that target. Consider eligible accounts at their first
servicing decision. The proposed comparison is the expected thirty-day net
cash recovery under two specified contact rules, including contact costs.
Follow-up ends at a common horizon; sale, transfer, default and missing
recovery need distinct definitions. A default state is absorbing only if
the economic outcome and action rules make it so. Evaluation of a fixed
rule, learning a rule, and constructing an interval for the value of a
selected rule require different statistical arguments.

Let $H_k$ contain the baseline, previous delivered actions and information
available before decision $k$. An event timestamp is insufficient: an event
on day seven that arrives on day nine is unavailable for a day-eight rule.
The assistant can propose this history and explain exclusions, but the
decision clock must be enforced by deterministic validation.
Let $A_k=g_k(H_k)$ be a fixed deterministic action rule, and let $K$ be the
last decision index. With complete observation, a discrete representation is
\begin{equation}
 V(g)=
 \sum_{\bar l}
 m\{h_K,g_K(h_K)\}p(l_0)
 \prod_{k=0}^{K-1}
 p\{l_{k+1}\mid h_k,g_k(h_k)\}.
 \label{eq:exec16-policy-recursion}
\end{equation}
Here $m(h_K,a_K)$ is the conditional expected net outcome. Each history in
the sum contains the actions already generated by $g$, and each later action
must be recomputed from that generated history. Integrals replace sums for
continuous states. Consistency, sequential exchangeability, support for the
rule and a justified observation law give this expression its causal
meaning. A fitted transition model supplies a numerical approximation; it
cannot supply those premises. The earlier retention and augmented-score
development remains the appropriate baseline when outcomes disappear
\citep{hernan2020}.

\subsection{What the direct collaboration literature establishes}

Zhang and colleagues' CausalCollab is a direct precedent for causal analysis
of sequential human--language-model interaction
\citep{exec13_census_collaboration}. Human refinements are treatments,
model replies are evolving states, and an evaluated final text is the
outcome. The paper recognizes that a particular text string is often an
inappropriate intervention across different contexts. It instead defines an
Incremental Stylistic Effect, or ISE, through a lower-dimensional function
of the refinement and its history. Its equation~(2) averages over text
versions conditional on the specified style and history. Its equation~(1)
then takes a limiting difference for an infinitesimal style shift.
Theorem~1 identifies that target under the paper's sequential
exchangeability and positivity conditions.

This construction is useful for the proposal because it makes the
intervention's versions explicit. It does not establish that any
embedding has sufficient causal information, that every desired style is
supported in every banking history, or that an arbitrary new generator
implements the same conditional version mixture.
Knowing that a text can be rewritten more formally does not demonstrate
that the required intervention has historical support among accounts
subject to the same eligibility, communication and distress conditions.
Likewise, a recorded conversation can omit a human's private judgement,
an earlier unlogged contact or information from another channel.
An LLM's description of the history cannot establish sequential
exchangeability in their presence.

Algorithm~1 proposes learning style representations with a conditional
variational autoencoder, fitting an outcome model and evaluating the
functional with Monte Carlo calculations.
The empirical evaluation uses rebuilt interaction datasets, generated
counterfactual refinements and outcome labels, and a constructed
confounding split. Table~2 reports a meaningful but qualified tradeoff.
For CoAuthor at $\alpha=0.2$, no adjustment gives mean squared errors
$0.188$ on observational examples and $0.353$ on counterfactual examples;
G-E plus CVAE gives $0.216$ and $0.219$. The counterfactual error falls
while the observational error rises. These are averages over three seeds
with reported standard deviations, not causal interval coverage
measurements or a uniform superiority result. PCA and CVAE yield similar
counterfactual errors in the table, so the method name does not by itself
justify extra complexity. The text labels and counterfactual generation
procedures define the experiment's truth construction; they do not
constitute observed customer cash outcomes
\citep{exec13_census_collaboration}.

The authors explicitly leave optimal collaboration strategies to future
work. Their ethical discussion also says that the rebuilt datasets are for
research and cannot be redistributed. Consequently, the public code and
dataset links do not establish a permitted product training corpus.
This paper supports a conditional text-intervention target and a promising
research example. It supplies neither a generally qualified longitudinal
policy learner nor evidence of bank deployment.

\subsection{A style intervention needs a version distribution}

For a one-step illustration, let $S=s$ denote a style and $T$ the actual
message. A proposed intervention kernel $q(t\mid s,h)$ specifies which
version is delivered in history $h$. Writing the version-specific
conditional outcome as $m(t,h)$, its value takes the form
\[
 V_q(s)=\sum_h p(h)\sum_t m(t,h)q(t\mid s,h).
\]
This is an authored simplification used to explain the deployment contract;
it does not replace the paper's longitudinal definition.
If two messages have the same style but different consequences,
changing $q$ changes the intervention even when $s$ is unchanged.
Suppose their stipulated gross values are 120 and 60 dollars and either
message costs eight dollars. A historical mixture assigning probability
$3/4$ to the first version has value
\[
 (3/4)120+(1/4)60-8=97.
\]
A generator assigning probability $1/4$ to that version has value 67.
The thirty-dollar difference comes entirely from the version mixture.
The example is reproducible arithmetic, not a measured language effect
\citep{exec16_longpolicy_oracle}.

There are several substantive remedies. The bank can freeze a finite
template set and its randomized assignment law, retain the complete
message as part of the intervention, or define and validate a
history-dependent generation kernel. An assumption that outcomes depend
on style alone would also remove version dependence in this example,
but it is a strong scientific restriction, not a benefit obtained merely
by compressing text. Changing a model, prompt, retrieval source or
decoding rule can change the kernel. A product proposal should therefore
bind those identities to the policy comparison and explain how a later
version would be reevaluated.

Locality creates another boundary. Suppose a stipulated gross response
curve is $v(s)=100+20s-40s^2$. At $s=0$, its derivative is positive:
$v'(0)=20$ dollars per style unit. Moving to $s=1$ nevertheless changes
gross value by $-20$; after an eight-dollar deployment cost, the finite
net change is $-28$.
An infinitesimal effect cannot be reported as the gain from that finite
change without an additional response and cost argument
\citep{exec16_longpolicy_oracle}. For banking use, the assistant should
distinguish a local diagnostic from a decision over feasible message
versions and doses. It should expose that difference before recommending
an action.

\subsection{The published code's actual statistical path}

The public CausalCollab repository at the inspected June 2024 revision
provides a research script with a useful, bounded implementation path
\citep{exec16_longpolicy_code}. Its caller dispatches to a logistic outcome
model, optional treatment representations, a learned last-state mean and
generated-state predictions. The default action count is two.
In \texttt{cal\_g\_estimate}, the code samples baseline embeddings from a
training bank, retains the supplied action or style embeddings, and
replaces the last state with generated draws before predicting the
outcome. Earlier nonbaseline states are retained from the input.
There is no call in this path that recomputes the next action as a function
of each newly generated history. This is a precise limitation on reuse
as an adaptive servicing evaluator; it is not evidence that the authors
claimed to implement such a banking evaluator.

The last-state sampler uses a fitted conditional mean and Normal noise
with default standard deviation one. That choice is a model for the
embedding distribution, not a universal transition or residual law.
The CVAE path sets its KL weight to zero; the inherited loss
backpropagates reconstruction error plus that weight times the KL term.
The inspected encoder and decoder are linear, with stochastic latent
sampling. Thus an adapter cannot infer a particular regularized latent
model merely from the label CVAE. It must bind the actual objective,
representation and sampling configuration.
The output path writes observed and constructed-counterfactual class
probabilities, the columns \texttt{gE0} and \texttt{gE1}, and MSE
comparisons. The inspected caller-to-output path does not expose an
explicit $\delta\to0$ derivative calculation or a confidence interval.
This is a path-specific source finding; no claim is made about unseen
external experiments or a package-wide impossibility
\citep{exec16_longpolicy_code}.

Several engineering conditions follow directly. The README's example
spells the representation option \texttt{CVAE}, while dispatch uses
\texttt{cVAE}; a proposed wrapper should validate the enumeration before
execution. Importing the utility module requests unrevisioned
\texttt{bert-base-uncased} weights and unconditionally moves the feature
model to CUDA after an earlier CPU branch.
An integration would need explicit device selection, pinned model
artifacts and an intentional acquisition boundary.
Shuffled row folds do not group by the participant or topic IDs retained
by the loader. If participants, households or calendar shocks span rows,
that split does not support an independence claim.
The appropriate grouping and time holdout must follow the deployment
population; their impact was not measured by this source inspection.

The nontruncated repository tree contains no root license, dependency
lock, package setup or test suite. These observations leave reuse rights,
installation closure and maintenance responsibility unresolved.
Twenty-two original blobs, including the README and selected caller,
loader, loss and output paths, were verified against the pinned tree.
Other dataset loaders and installed dependency behavior remain outside
this inspection. No source import or upstream experiment was executed.
The proposed decision is reference use pending rights and prototype
qualification. An independent implementation of the mathematical
baseline is a viable alternative, with its own attribution, review and
maintenance obligations; public visibility alone establishes neither
copying nor redistribution permission.

\subsection{Recomputing actions changes the banking answer}

The servicing law makes the missing recursion concrete. Hold the initial
action at $A_0=1$, so $P(L_1=1)=2/5$, and use the already specified
net outcome
\[
 Y=102-40L_1+(30L_1-8)A_1.
\]
The adaptive rule $A_1=L_1$ yields 102 in the low-distress state and 84
in the high-distress state, giving $V=94.8$.
Suppose a hypothetical converter instead draws an action with the correct
marginal probability $2/5$, independently of the newly generated distress
state. Its value is
\[
 102-40(2/5)+(2/5)\{30(2/5)-8\}=87.6.
\]
The loss of feedback changes value by 7.2 dollars per account.
Relative to making no second contact, whose value is 86 with this same
first action, the adaptive gain is 8.8 while the independent-action gain
is only 1.6. Matching action frequencies does not preserve a dynamic
intervention. Exact enumeration and separate backward and factorized
calculations agree on these quantities
\citep{exec16_longpolicy_oracle}.

This is a project-authored counterexample to the hypothetical conversion,
not an observed output of CausalCollab. Its purpose is to define an
acceptance case. A proposed adapter must regenerate $H_1$, call $g_1$
on that history, retain the resulting action, and propagate the complete
trajectory to the outcome. For more stages, it must repeat those steps.
It should agree with exact enumeration on small finite laws and with a
separately implemented backward recursion. A stochastic policy also needs
its conditional action probabilities and controlled random draws.
Padding a fixed vector of actions or reusing factual future actions is
insufficient.

Support may fail before any numerical discrepancy appears. If each of
four required actions has conditional logging probability $0.1$ along
the policy history, the complete matching path has probability $0.0001$.
For 10,000 independent trajectories, the expected matching count is one
and the probability of none is approximately $0.367861$.
Adding independent retention probability $0.8$ at each step gives an
expected complete matching count of $0.4096$ and a zero-count probability
of approximately $0.663910$
\citep{exec16_longpolicy_oracle}. These are planning calculations under
stipulated probabilities, not an estimate of a particular method's
variance. Outcome-model structure may permit borrowing across histories;
that adds assumptions and does not create missing randomized evidence.
Restricting the rule, changing the target population, collecting more
informative data or proposing a limited experiment are substantive
responses. Silently clipping large weights can change the statistical
object without resolving the original support question.

\subsection{Related frontiers and the limits of transfer}

The medical causal co-pilot roadmap by Alaa and colleagues helps organize
eligibility, treatment, follow-up, assumptions and accountable human
review. It explicitly discusses time-varying exposure and the danger of
convincing but wrong estimands; illustrative responses are not a
qualified sequential estimator
\citep{v12_la_2407_19118}.
The OCI workflow supplies a more concrete precedent for reproducible
analysis and human supervision, including point-treatment designs and
multiple simultaneous treatments through partially linear models.
Its simultaneous-treatment extension does not supply the evolving
history recursion required here
\citep{v12_la_2607_22443}.
These sources motivate an assistant that records alternatives, identifies
missing premises and hands a reviewed target to deterministic software.

Lampinen and colleagues show that active causal strategies can be
learned from passive training examples in their task environments.
Their Appendix~B.2 finds active intervention during evaluation critical
in the examined ablations. The positive learning result therefore does
not imply that observational banking logs alone settle the consequences
of a new policy
\citep{v12_la_2305_16183}.
The Model Discovery Agent distinguishes open-loop input sequences from
closed-loop policies and states that its experiments use open-loop
policies. Candidate mechanistic models and experiment proposals are
relevant assistance, but forecasting or approximate model evidence does
not independently identify an adaptive policy value
\citep{v12_la_2608_09696}.
Shi and colleagues' two-wave discovery result adds another useful
contrast: its linear non-Gaussian and higher-cumulant conditions address
path discovery, with nonzero-cumulant restrictions and power limitations.
They neither establish a longitudinal language-model estimator nor turn
edge detection into a consistently estimated monetary effect
\citep{exec15_latent_shi}.

The appropriate formal alternatives remain longitudinal g-computation,
weighting and recursively augmented or targeted estimation under their
specific assumptions. The inspected R \texttt{ltmle} interface is a
closer candidate for a restricted binary-treatment/censoring mean than
a text-collaboration demonstration, but it also requires correct node
order, dynamic-rule semantics, outcome transformations and observation
laws \citep{sw_ltmle}. None of these choices removes hidden confounding
by adding an LLM. A representation that predicts well may still omit
information needed for assignment or future response. Repeated model
answers, Monte Carlo variation and variation across fitted causal
models are different uncertainty sources; they cannot be exchanged for
sampling uncertainty or identification bounds.

\subsection{A bounded integration and public-evidence program}

The proposed CLI and MCP entrypoints should accept the same reviewed
scientific request. It binds population, unit, horizon, policy and
comparator; event, availability and decision clocks; outcome units and
costs; treatment and retention laws; and the graph or structural
assumptions used for identification. A text-generating intervention also
binds its templates or generation kernel, model artifacts and
representation. The policy is versioned independently of the
estimation code. A request that changes the policy during evaluation
requires a new analysis scope.

The response should distinguish an invalid chronology, an unsupported
backend, unavailable identifying information, failed support or
diagnostics, and a successful calculation under stated assumptions.
It should carry the admitted target, supported histories, numerical
error estimate where justified, statistical uncertainty method, source
identities and observed tool outcome. A JSON success field cannot
certify exchangeability.
An assistant may draft the request, explain refusals and compare remedies.
The proposed product should require an accountable review of its
scientific premises before estimation; it should not infer approval
from fluent prose or automatically send customer contacts.

The next affordable prototype has two complementary evidence tracks.
Finite generated mechanisms can expose known target and numerical
errors: feedback-aware versus fixed-action rollout, missing retention,
unsupported histories, changed text-version mixtures, late-arriving
states and local effects presented as finite gains. Exact enumeration
and independently coded backward recursion supply references. A
deterministic analyst workflow and a scripted assistant using the same
facts supply useful initial comparators. Where a genuine public
longitudinal dataset has suitable documentation and permissible use, it
can test ingestion, clock construction, diagnostics and reproduction
of a declared analysis. Its observed outcomes do not reveal both policy
counterfactuals. The rebuilt collaboration datasets should not be
redistributed on the strength of their public links.

The public-data survey supplies concrete choices. UCI's credit-card
default dataset contains monthly payment and bill histories with
documented CC BY 4.0 terms. It can supply covariate and timing exercises,
but a customer payment is not an assigned bank contact and the file
does not disclose the full servicing action law
\citep{pd15_uci_credit}. Explicitly generated actions and responses on
such covariates would constitute a semi-synthetic study.
Freddie Mac and Fannie Mae performance histories add realistic loan
and termination clocks, subject to registration and use restrictions;
their modification records do not supply all feasible actions,
private hardship information or known assignment probabilities
\citep{pd15_freddie,pd15_fannie}. They are conditional realism
candidates. The fully generated route remains usable while access and
rights are resolved.

Prototype acceptance should bind each expected target, unit, refusal
and numerical tolerance before running the comparison. Every attempted
case remains in the denominator, including missing evidence and tool
failures. A later live-model comparison needs frozen models, prompts,
tool budgets, held-out task mechanisms, a stopping rule and
prespecified uncertainty for its error and review-effort measures.
Separate policy learning from value evaluation; preserve account or
household groups and an appropriate future calendar holdout.
The exact arithmetic here executes reference calculations only. It
does not execute that prototype, a live-model trial or an upstream
package benchmark \citep{exec16_longpolicy_design}.

Costs include analyst formulation and review, rights and environment
resolution, embedding inference, nuisance fitting, trajectory
simulation, diagnostics, failed attempts and maintenance after a
model or policy revision. The inspected nested sampling path uses
50 baseline draws and 50 inner repetitions, amounting to 2,500
outcome-prediction rows per evaluated case. Shared fitted models and
baseline draws mean this is not 2,500 independent observations, nor is
the loop count a runtime measurement
\citep{exec16_longpolicy_code,exec16_longpolicy_oracle}.
A fair comparison measures the full workflow for the same policy,
information and accepted uncertainty. Real prices, realized recovery
and institutional adoption remain later inputs.

The engineering repairs have known routes: explicit clocks and policy
recursion, model pinning, validated options, grouped evaluation,
typed failures and exact behavioral references. Missing reuse terms
and uninspected originals require evidence acquisition or a permitted
alternative. General representation sufficiency, latent longitudinal
identification, valid uncertainty after policy selection, and useful
precision in poor-support regimes remain research obligations.
The proposal can justify pursuing those questions with public and
generated evidence while leaving their answers open
\citep{exec16_longpolicy_design}.
\label{exec16:longitudinal-policy-end}


% SOURCE docs/release/2026-10-29/source/docs/writing-repair/2026-10-29/src/41d61ab9a3845f5b.tex CHAR 0
\hypertarget{agent-architectures-must-preserve-scientific-uncertainty}{%
\subsection{Agent architectures must preserve scientific
uncertainty}\label{agent-architectures-must-preserve-scientific-uncertainty}}

\label{exec13:census-agents}

Several systems now make causal analysis a sequence of tool-supported
decisions. Their strongest contribution is often a concrete
representation of the intermediate work: a dataset, a graph, an
estimand, a proof statement or a failed tool call. Comparing these
systems requires examining what each intermediate object means and which
transformations are allowed. A graph stored in agent memory can prevent
a variable-name error while still carrying an unsupported orientation
into every later calculation.

\hypertarget{causal-agent-useful-tool-boundaries-and-unequal-baselines}{%
\subsubsection{Causal Agent: useful tool boundaries and unequal
baselines}\label{causal-agent-useful-tool-boundaries-and-unequal-baselines}}

Han and colleagues store causal graph objects in a dictionary that
agents address by name. This connects independence tests, graph
questions and effect calculations without requiring each operation to
reconstruct the graph from prose. Their CausalTQA examples combine small
generated nonlinear additive Gaussian mechanisms with medical and
marketing vocabulary. The vocabulary tests a language interface; it does
not turn the generated mechanism into a medical or marketing
observation. \cite{exec13_census_causal_agent}

The paper documents informative failures: requesting conditional rather
than marginal independence, passing the wrong variable name, using the
wrong graph operation, and building unnecessarily long tool chains.
Typed variable references and explicit tool signatures are plausible
remedies. The main code-generation baseline, however, executes once,
whereas the agent uses specialized tools and iterative reasoning. Its
measured advantage therefore combines tool design, supplied knowledge
and opportunities for repair. An evaluation intended to isolate
orchestration should give a deterministic workflow and a simpler agent
the same tool implementations and permitted repair budget.

The reported average cost of 3.37 seconds, 2.51 iterations and 5,308
tokens comes from 100 sampled variable, edge and graph questions with a
maximum iteration count of 15. It does not establish the cost of the
full effect workflow, rare failures or deployment under a different
model and platform. An average well below a limit is also insufficient
evidence to remove that limit. The prototype should retain explicit
budgets and report tails, timeouts, failed requests and review effort.
The source's effect-answer tolerance and benchmark oracle require code
inspection before its reported 157 correct answers out of 160 become a
reproduction target. \cite{exec13_census_causal_agent}

\hypertarget{orca-shared-state-without-silently-choosing-a-causal-world}{%
\subsubsection{ORCA: shared state without silently choosing a causal
world}\label{orca-shared-state-without-silently-choosing-a-causal-world}}

ORCA organizes analysis of relational databases around dataset
exploration, discovery and inference. It records table choices, joins,
variable roles, candidate graphs and estimation configurations, with
opportunities for human review. Its REEF benchmark uses a specified
structural causal model inside a relational e-commerce database. This is
directly relevant to a banking proposal: a mistaken join or aggregation
can change the unit of analysis before a causal estimator is called.
\cite{exec13_census_orca}

The important limitation lies between graph discovery and inference. The
paper describes a consensus PAG that retains uncertain endpoints,
followed by a separate ``actionable'\,' DAG. Uncertain directions are
resolved using the highest-ranked candidate, cycle-safe orientation,
reversal or edge removal. Acyclicity does not establish that the
selected orientation is true. An effect computed from this one DAG is
conditional on choices added during that conversion, and cannot
automatically represent all models encoded by the uncertainty-preserving
graph.

For a servicing decision, suppose distress and contact intensity remain
ambiguously related. Directing the edge one way can make distress an
adjustment variable; directing it the other can make it a treatment
descendant. The proposed interface should retain the original graph
object and give every added orientation an explicit assumption record.
It should compare identification across compatible models, return an
effect set where appropriate, or state that a scalar target is not
identified. Ranking an orientation for display is not authorization to
treat it as scientific fact. This is a proposed restriction on reuse,
not behavior shown to exist in ORCA's implementation.

The reported user study also requires careful interpretation. Twelve
participants used ORCA and a baseline with tools, step gates and minimal
state in counterbalanced order. Table 1 gives rounded completion rates
of 0.83 and 0.75. The abstract's 42-percentage-point completion gain is
not consistent with that table; it should not be repeated as a verified
result. Appendix E defines completion as producing an ATE regardless of
accuracy, and discovery success as including both treatment and outcome
nodes. Neither definition establishes correct discovery.
\cite{exec13_census_orca}

The paper's ``CI coverage'\,' description counts an estimated ATE
falling within an interval around a ground-truth effect. That wording
differs from the usual repeated-sampling event in which an estimated
interval contains a fixed true effect. Its exact oracle and
implementation require clarification. Likewise, the full-data test for
the ATE-error condition effect has \(p=.321\), whereas the analysis
excluding large configuration-error outliers reports \(p=.020\) under
OLS. Configuration failures are central outcomes for a workflow product.
They should remain in the primary comparison, with conditional numerical
accuracy reported separately.

\hypertarget{causalpulse-operational-success-has-a-different-denominator}{%
\subsubsection{CausalPulse: operational success has a different
denominator}\label{causalpulse-operational-success-has-a-different-denominator}}

CausalPulse supplies another concrete design: agents coordinate semantic
descriptions, domain rules and causal tools through a utility layer that
includes MCP. Such interfaces can make tools accessible and operations
traceable. The inspected manufacturing study nevertheless distinguishes
successful operation from successful root-cause attribution. Its PSE
example reports Hits@1 of .44, precision .46, recall .40 and F1 .39
despite much higher operational success. Tool availability and
completion are necessary engineering outcomes, but substituting them for
causal correctness would misstate the evidence. \cite{exec13_pulse}

This distinction determines the proposed monitoring interface. It should
record whether the request was parsed, whether the analysis was
admissible, whether the tools completed, whether the scientific target
was answered, and whether the decision was useful under the declared
loss. These events need separate denominators. A scientific refusal can
be a correct answer, while a complete report can be wrong. Banking
acceptance would additionally need appropriate target-bank evidence; the
public prototype can already test these distinctions with generated
known-truth cases.

\hypertarget{causalsmith-a-checked-proof-still-needs-a-checked-statement}{%
\subsubsection{CausalSmith: a checked proof still needs a checked
statement}\label{causalsmith-a-checked-proof-still-needs-a-checked-statement}}

CausalSmith combines a Lean library with proposal, formalization, proof
and presentation stages. Its paper reports 8,179 declarations and a
catalog of 144 runs: 14 accepted, 59 downgraded and 71 failed. Those
counts include reruns as separate entries and exclude a legacy track.
They describe one campaign with changing model versions, not independent
draws from a fixed research-task population. The authors explicitly
acknowledge that research quality ratings lack independent human
validation and that the error rate of their LLM statement-match audit is
unknown. \cite{exec13_census_causalsmith}

The source's most useful distinction is between checking a proof against
a formal statement and checking whether that statement expresses the
intended claim. A compiled artifact can assume a hard step as an axiom,
weaken a conclusion, strengthen an assumption or make a predicate
vacuous. The paper describes mechanical scans and dependency checks
alongside language-model statement review. These are complementary
checks with different assurances. The latter does not acquire a formal
guarantee merely by sitting next to a proof checker.

For a proposed bank-effect theorem, the statement record should expose
the unit, intervention, estimand, time ordering, conditioning
information, positivity and dependence assumptions. The proof
environment should bind the toolchain, imported definitions and axioms.
A separate review should compare the English claim with that exact
formal statement and test whether its assumptions are satisfiable in the
intended model. A theorem conditional on observed confounders cannot be
presented as solving unmeasured confounding in operational data. Formal
correctness, semantic fidelity and scientific relevance each need their
own evidence.

The near-term reuse choice is therefore selective. Retain typed graph
memory, explicit intermediate state, preserved failures, and formal
statement/dependency records as design candidates. Defer executable
adoption of these newly inspected systems until exact code,
dependencies, licenses and relevant numerical or proof paths have been
reviewed. This permits a concrete prototype architecture now without
claiming that paper descriptions have already been integrated or
independently reproduced.

\label{exec13:census-agents-end}


% SOURCE docs/release/2026-10-29/source/docs/writing-repair/2026-10-29/src/ded75c1791455c8b.tex CHAR 0
\section{Choosing an agent architecture for a servicing decision}
\label{exec13:agent-architecture}

A bank considering assisted causal analysis faces an architectural choice
before it faces a model leaderboard. It can give a language model a broad
notebook environment, surround a fixed statistical pipeline with conversational
assistance, or ask several specialist agents to negotiate an analysis. Each
choice changes which mistakes are easy to detect, what must be reviewed by a
person, and how much an apparently completed answer costs. The useful comparison
is therefore between complete analytical paths with declared scientific
obligations. Counting agent roles or successful tool calls gives little guidance
about whether a servicing decision has been answered.

Consider a manager asking: ``Should we offer an early restructuring review to
customers entering arrears?'' Let the unit be an eligible account at its first
qualifying arrears date, $X$ the information available immediately before the
offer, $A\in\{0,1\}$ the offer, and $Y$ the discounted net cash recovery over a
fixed following horizon. The offer is distinct from acceptance and from the
eventual restructuring terms. For a specified eligible population $P$, an
initial target is
\begin{equation}
 \tau_P=E_P\{Y(1)-Y(0)\}.
 \label{exec13:architecture-target}
\end{equation}
If $Y$ excludes the incremental cost $c$ of making the offer, an expected-value
decision compares $\tau_P$ with $c$. If costs are already inside $Y$, subtracting
them again is wrong. An agent should make this definition explicit before
requesting an estimator. Neither an offer-acceptance prediction nor the effect
among customers who accepted is automatically an answer to this question.

The information clock is equally consequential. A field can carry an event
date before the offer but be entered into the database after it. The admissible
record is the version available at the decision time, with the extraction and
revision rule specified. A distress assessment obtained during the offer call
cannot silently become a pretreatment covariate. Multiple accounts held by one
customer, branch policies and shared collection teams require an explicit
sampling and interference analysis; calling each row an independent customer
does not settle these issues. The simple point-treatment example below assumes
one eligible unit, no interference and a declared independent-unit law. The
longitudinal construction earlier in this chapter is needed when subsequent
contacts and evolving distress become part of the intervention.

\subsection{What the competing architectures actually contribute}

The primary literature supports several distinct proposals. Zhu and colleagues'
CausalAgent routes processing, structure learning and reporting through agents
and MCP tools. Its four-page demonstration uses the Sachs dataset; quantitative
effect estimation and counterfactual modules remain among its stated future
directions. Its conversational interface is evidence of an interaction design,
not a demonstrated solution to the servicing target in
\eqref{exec13:architecture-target}. Han and colleagues' separately authored
Causal Agent instead illustrates named graph objects in tool memory and a
generated question benchmark. Keeping these two systems distinct prevents a
result from one from being attributed to the other's implementation
\cite{v12_la_2602_11527,exec13_census_causal_agent}.

The actor/critic design of oci-agent addresses a closer statistical problem:
point-in-time observational effects under supplied unconfoundedness. It exposes
specifications, notebooks and diagnostics to a human principal. This can reduce
the effort of reproducing an analysis or finding a failed overlap diagnostic.
The critical question is whether a proposed revision preserves the target.
Removing unsupported parts of the population may make an estimate more stable
while answering a different banking question. That is a substantive decision
which an attractive report must not hide. The paper's benchmark results remain
reported source findings, and are not reproduced live-model results of this
proposal \cite{v12_la_2607_22443}.

The causal-learn+ position takes a stricter approach. Agents assist data
inspection, method explanation and interpretation around a protected discovery
core. Their proposed edges, orientations or constraints do not automatically
become scientific evidence. This reduces one path by which language-model
beliefs enter a graph, but leaves assumption elicitation, variable construction
and statistical method choice as important tasks. Its call for future
benchmarks is consistent with treating the architecture as a candidate for
evaluation rather than an established performance advantage
\cite{v12_la_2606_23608}.

ORCA emphasizes shared relational state: joins, variable roles, graphs and
estimation configurations. This is especially relevant when banking data are
spread over account, customer, contact and transaction tables. However, the
paper's conversion of a consensus PAG into an actionable DAG introduces
orientation, reversal or deletion decisions. An analysis conducted on that DAG
inherits these extra choices. The architecture should expose them as additions
to the assumptions, and preserve the original uncertain graph. The numerical
comparison is also narrower than its headline suggests: the inspected table
reports completion rates of $.83$ and $.75$, whereas the abstract claims a
42-percentage-point gain. Completion is defined as producing an ATE regardless
of its accuracy. These are reasons to preserve the source discrepancy and to
design a better denominator, not to invent a corrected experiment
\cite{exec13_census_orca}.

Two further architectures broaden the comparison. CausalSmith makes a formal
statement, proof dependencies and compiler results explicit; its own limitations
separate proof checking from unvalidated model judgments about statement
faithfulness and research quality. Model Discovery Agent proposes executable
mechanisms, then uses numerical fitting and experiment design to compare them.
The latter's model-disagreement criterion approximates information gain within
a changing candidate set. Its reported predictive tests do not establish
causal identification, and its illustrated experiments use open-loop policies.
Neither formal proof search nor mechanism proposal alone replaces the
observational assumptions needed for the servicing target
\cite{exec13_census_causalsmith,v12_la_2608_09696}.

\begin{longtable}{@{}P{.19\textwidth}P{.35\textwidth}P{.39\textwidth}@{}}
\caption{Architecture choices for a bounded servicing prototype. These are
proposal decisions based on inspected sources.}\label{exec13:architecture-choices}\\
\toprule
Design & Useful intermediate object & Proposed reuse boundary\\\midrule
\endfirsthead
\toprule Design & Useful intermediate object & Proposed reuse boundary\\\midrule
\endhead
Conversational causal tools & Named variables, graph objects and explicit
tool arguments & Adapt interaction ideas; independently check graph encoding,
tool scope and malformed-result handling.\\
Actor/critic notebooks & Analysis specification, code, diagnostics and
criticism & Adapt review artifacts; restrict execution to a frozen supported
analysis and require a new target declaration after trimming.\\
Protected discovery core & Distinct suggestions, decisions, parameters and
graph outputs & Prefer for an initial discovery assistant; retain graph
ambiguity and test the advice surrounding the core.\\
Shared relational state & Joins, populations, candidate graphs and estimators
& Adopt the concept of typed intermediate state; keep every consequential
join and orientation inside the scientific review.\\
Formal proof workflow & Exact statement, definitions, imports and proof
& Use for selected derivations; review semantic fidelity and satisfiability
separately from compilation.\\
Mechanism proposal loop & Executable mechanisms, fitted weights and experiment
candidates & Consider for a later controlled simulator; inspect model support,
action coverage and approximation error before causal interpretation.\\
\bottomrule
\end{longtable}

These alternatives suggest a restrained first prototype: one application-owned
controller, a bounded set of numerical and graph operations, and optional
language-model assistance for drafting and criticism. Multiple personas can be
an experimental arm. Their added value should be measured against the same
tools, information and total repair budget available to a simpler workflow.
This choice concerns the proposed product; it does not claim that the reviewed
external frameworks have been installed or integrated.

\subsection{A target can drift while every tool succeeds}

Suppose consistency, conditional exchangeability
$(Y(0),Y(1))\perp A\mid X$, positivity and the stated sampling law hold.
Write $e(x)=P(A=1\mid X=x)$ and $m_a(x)=E(Y\mid A=a,X=x)$.
The familiar augmented inverse-probability score is
\begin{align}
 \phi(O;m,e)={}&m_1(X)-m_0(X)
   +\frac{A\{Y-m_1(X)\}}{e(X)}\nonumber\\
 &-\frac{(1-A)\{Y-m_0(X)\}}{1-e(X)}.
 \label{exec13:architecture-score}
\end{align}
At the true nuisance functions, conditioning on $X$ removes both residual
terms, leaving $E\{\phi\mid X\}=E\{Y(1)-Y(0)\mid X\}$.
Consequently, averaging over $P$ gives $\tau_P$. Estimated nuisance functions,
cross-fitting, dependence and uncertainty require the additional arguments
developed elsewhere in this monograph. Calling this score does not establish
exchangeability or those rate conditions \cite{hernan2020}.

Now let an overlap diagnostic cause the workflow to keep only
$K=\{x:.1\leq e(x)\leq .9\}$. A uniform average of the score among retained
units targets
\begin{equation}
 \tau_K=\frac{E[1_K(X)\{m_1(X)-m_0(X)\}]}{P(X\in K)}.
 \label{exec13:architecture-trim}
\end{equation}
By contrast, define the overlap-weighted target explicitly as
\begin{equation}
 \tau_{\mathrm{ov}}=
 \frac{E[e(X)\{1-e(X)\}\{m_1(X)-m_0(X)\}]}
      {E[e(X)\{1-e(X)\}]}.
 \label{exec13:architecture-overlap}
\end{equation}
Equations \eqref{exec13:architecture-trim} and
\eqref{exec13:architecture-overlap} integrate against different measures.
They coincide only under additional restrictions, such as a constant
conditional effect; their labels cannot make them equal. If $K$ is estimated
from the same data, uncertainty about the selection rule is another issue.
The present calculation uses known propensities and a fixed rule so that this
additional difficulty does not obscure the target distinction.

In the inspected oci-agent notebook, the branch labeled ATO calls a propensity
trimming function and then applies a uniform AIPW summary to the retained
pseudo-outcomes. The serialization retains this label and exports retained
indices separately. This source path supports a trimmed-population
interpretation under the relevant assumptions; it does not implement the
weighting in \eqref{exec13:architecture-overlap}. The appropriate adaptation
therefore reads the implemented weights and retained population, not just the
display name \cite{v12_la_package_oci_agent}.

An exact generated example makes the financial consequence visible. Take
three strata with masses $(1/4,1/2,1/4)$, offer probabilities
$(.05,.5,.8)$ and conditional net-recovery effects $(-200,100,500)$ dollars.
Direct summation gives
\begin{equation}
 \tau_P=125,\qquad
 \tau_K=\frac{700}{3}\simeq233.33,\qquad
 \tau_{\mathrm{ov}}=\frac{48200}{283}\simeq170.32
 \quad\text{USD per account}.
 \label{exec13:architecture-target-numbers}
\end{equation}
Trimming retains three quarters of the original population. If the offer
cost is 180 dollars per eligible account and lies outside the recovery
outcome, the full-population expected net gain is $-55$ dollars. Substituting
the trimmed effect into that full-population calculation would reverse the
sign. A policy confined to the retained stratum is a different decision and
needs its own eligibility rule. The numbers are authored potential outcomes,
not observations about any bank. A separate 400-unit finite-population
calculation with exact nuisances reproduces all three weighted averages
\cite{exec13_architecture_exact}.

The same problem occurs in graph handoffs. Under one declared DAG, a baseline
distress measure can be a common cause of offer and recovery and belong in an
adjustment set. Under a different timing definition it can be a descendant of
the offer, so conditioning on it can remove part of the total effect. An
architecture should attach the variable definition, information time and graph
assumption to the adjustment result. Caching by column names alone cannot
preserve that meaning.

\subsection{What repeated critics can and cannot repair}

A critic is useful when it brings a distinct check: reconstructing weights,
reading a source passage, evaluating a graph criterion or reproducing a
numerical baseline. Merely repeating the same prompt with the same evidence
does not create independent scientific information. To quantify this
distinction, suppose each of three critics misses an error with marginal
probability $q=.1$. If their errors are independent, majority failure has
probability
\begin{equation}
 q^3+3q^2(1-q)=.028.
 \label{exec13:architecture-majority}
\end{equation}
Perfectly shared errors instead leave the failure probability at $.1$.
Consider an intermediate stipulated law: with probability $\gamma=.05$,
all three fail from a common cause; otherwise they fail independently with
probability $r=1/19$. Each still has marginal error
$\gamma+(1-\gamma)r=.1$, but majority failure becomes
\begin{equation}
 \gamma+(1-\gamma)\{r^3+3r^2(1-r)\}\simeq .05762.
 \label{exec13:architecture-common-error}
\end{equation}
Thus the same individual error rate can accompany quite different benefits
from voting. Exact enumeration over all eight independent-error patterns
checks the calculation. None of these probabilities is an empirical estimate
of a language model or a deployed critic \cite{exec13_architecture_exact}.

There is also an information limit. Consider two linear Gaussian worlds with
mutually independent standard normal noises in each world:
\begin{align}
 \mathcal W_1:&\quad X=U,\quad Y=\tfrac12X+\sqrt{3/4}\,V,\nonumber\\
 \mathcal W_2:&\quad Y=V',\quad X=\tfrac12Y+\sqrt{3/4}\,U'.
 \label{exec13:architecture-worlds}
\end{align}
Both observed pairs have zero mean, unit marginal variances and covariance
$1/2$, hence the same joint distribution. The effect of increasing $X$ by
one is $1/2$ in $\mathcal W_1$ and zero in $\mathcal W_2$. Suppose an
adaptive agent system receives only samples from that common law, identical
documents in both worlds and randomness independent of the world. The first
transcript step then has the same law in both worlds. If this holds through
step $t$, the next tool or agent applies the same conditional transition to
an identically distributed history, so it holds at $t+1$. Induction gives the
same distribution for the final answer. An additional critic operating on
that transcript cannot distinguish these worlds.

This argument concerns the stated information set. Genuine temporal knowledge,
an intervention or a justified scientific restriction can break its premise.
Retrieval is valuable when it supplies such information with an inspectable
provenance; fluent agreement alone cannot do so. The appropriate output before
the added information is available can be an identified set, a conditional
answer or a request to clarify the decision model. It need not be an arbitrary
scalar produced to make the workflow look complete.

Repeated repair also changes evaluation. In a deliberately simple experiment,
suppose a scientifically invalid candidate independently has probability $.05$
of passing a fallible screen on each try. Allowing twenty tries and retaining
the first apparent pass gives probability
$1-.95^{20}\simeq.6415$ of at least one false pass. Without independence, the
union bound is only $\min(1,20\times.05)=1$. This example describes independent
screening attempts, not deterministic replay of the same input, and does not
assign a statistical significance level to real agent review. It explains why
the proposed evaluation must retain all attempts, give comparators the same
repair budget and separate development feedback from held-out evaluation
\cite{exec13_architecture_exact}.

\subsection{Typed handoffs and the scope of recovery}

For the servicing prototype, a useful analysis state contains the decision
contract, data identity and clock, graph or graph class, assumptions with their
source and status, target measure, method applicability, numerical artifacts,
uncertainty statement and request history. These objects should refer to one
another by immutable identity. A revised target or graph invalidates dependent
results even when the data bytes are unchanged. A successful statistical call
must not make an unresolved assumption disappear.

The controller can enforce local invariants. An estimation request must name a
target and population; every variable must exist in the supplied schema; a
numerical interval must have finite, ordered endpoints unless the result type
explicitly supports an unbounded set; and a result must bind the method and
input versions that produced it. A graph-identification operation can check a
formal condition relative to a supplied graph. These checks support meaningful
handoffs, but none proves that the graph or observation law describes the bank.
Scientific assumption status should therefore remain visible alongside
mechanical validity.

This distinction leads to several result types with different continuation
actions. Malformed input calls for correction of the request. A missing backend
calls for an implementation or environment remedy. A mathematical
nonidentification result calls for different information or a different target.
An unsupported method calls for a scoped alternative. A valid set-valued answer
can proceed to a decision calculation that respects the set. Collapsing all of
these into either an empty successful object or a single generic failure would
discard the reason that determines the next useful action.

The inspected CausalAgent server illustrates a concrete boundary to repair
before reuse. Its OLC handler refers to a function not imported or defined in
that module and catches the resulting exception into a JSON error string.
Separately, its output helper extracts text between braces and returns an
empty dictionary on a decoding failure. These are static source findings, not
a claim that this proposal executed the upstream server. A wrapper must
validate the scientific result and propagate an explicit failure before
downstream analysis; successful transport of an error string is not a learned
graph \cite{v12_la_package_causalagent}.

Recovery has a related distinction. A graph checkpoint preserves enough state
to resume an operation; it does not automatically make external effects occur
exactly once. The inspected LangGraph interrupt restarts its node from the
beginning. An application should therefore give an expensive computation a
stable request identity, record its durable outcome and check that identity on
replay. For a remote provider call interrupted before a receipt is saved, the
outcome can remain unknown unless the provider supports lookup or idempotency.
Blind resubmission cannot be advertised as exact-once execution
\cite{v12_la04_langgraph}.

Pydantic AI supplies typed tool boundaries and usage controls, but the inspected
price check cannot enforce a monetary limit when price is unavailable. Some
usage becomes known only after a response. A prototype requiring a hard budget
needs a known pricing contract, conservative reservation and explicit behavior
for unknown cost. AutoGen's inspected save/load methods preserve model context;
they do not constitute a journal of arbitrary external side effects. Its
maintenance declaration and separate code/documentation licenses also enter
the adoption decision. These observations favor a plain controller baseline
before adding a framework whose incremental benefit has not yet been measured
\cite{v12_la04_pydantic_ai,v12_la04_autogen}.

The existing local prototype provides narrower positive evidence. Six authored
offline cases preserve blocked, unsupported, ambiguous and completed outcomes,
including directed-revision recovery and a weak moment set. The saved replay
reproduces those outcomes using real local tools with no live model calls.
This supports selected semantics and recovery behavior. It does not test
whether a language model elicits the right servicing graph, reduces human
effort or behaves correctly on an external platform. The library, CLI and MCP
contracts should carry this same distinction when transporting the result
\cite{exec13_architecture_local}.

\subsection{Selecting the next experiment by useful output and total cost}

The minimal comparison has a deterministic analyst template, a single assisted
controller and, if justified, a controller with a critic. All receive the same
documents, permitted tool implementations and task budget. Generated cases
should include correctly identified effects, ambiguity requiring a set or
refusal, misleading variable clocks, a target change after trimming and a tool
failure. A held-out task split separates adaptation of prompts and templates
from evaluation. Numerical truth comes from the declared generator, while
semantic scoring needs an independently specified reference and ambiguity
rubric. A correct refusal is a successful answer to a nonidentified task;
unnecessary refusal on an identified task is a separate error.

Let $N$ count all eligible assigned requests and $U_i$ indicate that request
$i$ produced an admissible, correct and decision-appropriate answer under the
rubric. A useful workflow summary is
\begin{equation}
 \widehat C_{\mathrm{usable}}=
 \frac{\sum_{i=1}^N(C_i^{\mathrm{tools}}+wT_i^{\mathrm{review}}
                         +C_i^{\mathrm{other}})}{\sum_{i=1}^N U_i}.
 \label{exec13:architecture-cost}
\end{equation}
It is undefined when there is no usable output. Costs for failed requests stay
in the numerator. The statistic should accompany the usable fraction, error
categories, refusal behavior, latency distribution and conditional numerical
accuracy. Reporting only accuracy among completed estimates can reward a system
that abandons difficult requests or hides configuration errors.

For planning, suppose manual review costs 15 minutes per request at 120 dollars
per hour and yields a usable fraction $.9$. Its ratio is $30/.9=33.33$ dollars.
An assisted workflow costing 3 dollars in tools plus 6 review minutes, with
usable fraction $.8$, gives $15/.8=18.75$ dollars. At 20 review minutes the
same tool and quality assumptions give $43/.8=53.75$ dollars. Equating the
assisted ratio to the manual ratio yields a review-time threshold of
$71/6\simeq11.83$ minutes. These stipulated scenarios omit fixed build costs;
they show which quantities to measure, not that assistance saves money
\cite{exec13_architecture_exact}.

The next justified experiment is consequently a bounded public/generated task
study measuring semantic failures and total review effort, with matched tools
and frozen denominators. A live-model performance claim requires a recorded
backend revision, actual calls and measured usage. A literature and architecture
proposal can proceed without that claim. Authentic bank experiments become
worth considering when the prototype isolates an uncertainty whose resolution
could change the build or investment decision. The unresolved questions are
specific: how often assumptions are elicited correctly, whether criticism adds
independent information, how much adjudication is needed and whether the
resulting gain survives a change in task population. None is answered by adding
agent roles to a diagram.
\label{exec13:agent-architecture-end}


% SOURCE docs/release/2026-10-29/source/docs/writing-repair/2026-10-29/src/85372947f74831ad.tex CHAR 0
\section{Data and code agents: a fail-closed banking boundary}
\label{exec13:data-code-agents}

An agent that prepares data or code can save analyst time while silently changing a causal question. The proposed boundary starts with a typed decision contract and uses the agent for retrieval, clarification, code drafting and criticism around fixed operations. It does not execute arbitrary notebook or shell text, download mutable ontologies, or convert malformed output into an empty successful analysis. This is a proposal architecture supported by inspected papers, package source and local adverse fixtures; it is not a live-model performance result.

Suppose a bank asks whether an early arrears review should be offered to eligible accounts. The unit is the account at its first qualifying arrears date. The action $A$ is an offer made using data available immediately before that date, and $Y$ is discounted net recovery over 180 days. The target is $\tau_P=E_P[Y(1)-Y(0)]$ for the declared eligible population $P$, with the offer cost either inside $Y$ or subtracted once in the decision rule. A field's event date and its database availability date are separate: a hardship flag entered on day $11$ cannot be a day-$10$ covariate. The agent may list candidate fields and ask an analyst to resolve ownership, timing, interference and missingness; it cannot infer those facts from column names.

The non-agent baseline is a structured analyst interview, a versioned data dictionary and manually reviewed code. The agent-assisted path adds a typed specification, source-to-claim links, a graph candidate set, fixed operation requests and deterministic report rendering. Both paths receive the same information and the same numerical fixtures. The comparison should score target preservation, clock errors, unsupported certainty, unresolved questions, artifact completeness and analyst time. A shorter answer is not better if it changes $P$, $A$, $Y$ or the horizon.

\subsection{What the inspected frontier contributes}

CausalAgent describes processing, structure-learning and reporting agents routed through MCP tools and demonstrates a Sachs protein-signalling workflow. Its paper lists quantitative effect estimation and counterfactual modules as future work. It therefore supplies an interaction and routing idea, not evidence that a bank servicing effect was estimated. The inspected pinned server declares PC and OLC tools, but its OLC call is not imported or defined in the selected module. Its JSON helper extracts text between the first and last braces and returns an empty dictionary on decoding failure. The proposal retains a typed refusal or malformed-result state instead of accepting that empty object as success. \cite{exec13_data_causalagent,exec13_data_causalagent_code}

The human-augmenting OCI workflow is closer to this target. It places a human principal over specifications, templated notebooks, actor/critic review and overlap diagnostics, and reports public ACIC and LaLonde evaluations. Those are source-reported evaluations, not a live run here. Its pinned notebook labels a uniform AIPW mean after propensity trimming as ATO. With two equally frequent strata, propensities $0.2,0.5$ and effects $0,10$, the retained-population mean is $5$, while conventional overlap weighting gives $6.09756$. The label cannot change the estimand. Our adapter would retain retained indices, weights, target and trimming rule. \cite{exec13_data_oci,exec13_data_oci_code}

CausalDS is useful as a benchmark design because it separates scenes, observations, identification decisions, private structural ground truth and scoring artifacts. Its scorer treats abstention as correct only when the task is explicitly abstainable and the model intentionally returns a null-equivalent answer. It separately records interval width, coverage and interval score, with both uncapped and capped normalized scores. The source also has nullable identification states and branches where missing metadata can default toward an identifiable interpretation. A repair must guard missing metadata, finite interval endpoints, malformed answers and nonidentifiable targets before scoring. \cite{exec13_data_causalds,exec13_data_causalds_code}

OntoLearner provides modular ontology loading, typing, taxonomy and relation discovery. Its code can download ontology files from Hugging Face without a revision argument. A source-code commit therefore does not identify the ontology bytes used by an analysis. The bounded remedy is local, hash-bound ontology input with declared units and clocks; taxonomy membership and relation extraction remain semantic suggestions, not causal-edge truth or ownership evidence. \cite{exec13_data_ontolearner,exec13_data_ontolearner_code}

Causal discovery in the era of agents proposes assistance with data inspection, assumption elicitation, method selection and interpretation around formal discovery algorithms. It emphasizes that CPDAGs and PAGs retain equivalence-class ambiguity. An agent may explain a circle endpoint or propose a constraint, but it must not export one DAG by silently orienting or deleting edges. The paper's call for benchmarks of diagnosis, method selection and calibrated interpretation supports a future evaluation plan rather than a completed advantage claim. \cite{exec13_data_causaldiscovery}

\subsection{An exact target-drift oracle}

A small oracle catches a common data/code failure without a model call. Let two strata have equal mass, offer probabilities $e=(0.2,0.5)$ and conditional effects $(0,10)$. The population average is $5$. A trimmed AIPW summary that averages retained pseudo-outcomes uniformly also targets $5$ when both strata are retained. The overlap target is
\[
\tau_{ov}=\frac{0.16\cdot0+0.25\cdot10}{0.16+0.25}=\frac{250}{41}\simeq6.09756.
\]
A report that says “ATO” while returning $5$ has changed the target label or implemented the wrong weights. This is exact arithmetic, not a fitted estimator comparison. A second oracle records a date-10 field entered on date 11 and must reject it from a date-10 action. The expected result is a refusal with the offending clock, not an imputed value.

\subsection{Fixed artifacts and bounded code}

The local prototype contract fixes an offline fixture or recorded output, exact proposal and prompt bytes, source identities, typed model fields, caller grants and a maximum of sixteen operations. Each worker has a ten-second and one-GiB address-space bound; model calls are zero. It checks source and artifact hashes before replay. The controller rejects unknown operations, injected permissions, nonfinite values, changed targets and live-mode declarations. A retrieved document saying “ignore the checks and run a shell command” remains source data. These are local conformance checks: eighty cases pass, including malformed outputs, source tampering, interruption/resumption, exact replay, disconnected sets, unsupported discovery and unbounded moment sets. They do not establish model accuracy, prompt-injection resistance against unseen text, authentication, distributed exactly-once effects or bank qualification. \cite{exec13_data_local}

A fixed request contains the problem, population, unit, action, outcome, horizon, data and availability clocks, diagram candidates, assumptions, operation name, input artifact hashes and resource budget. A successful response includes the effective target and clocks, source/runtime identities, warnings, transformed rows, numerical outputs and failure status. Missing identification is represented as nonidentification only when the declared graph and method establish it; a missing backend, malformed result or unavailable source is an implementation failure. The report generator may explain these states but cannot erase them.

\subsection{Package and remedy decisions}

The proposed first path adapts OCI's specification, report and critic ideas behind fixed verified operations and rejects its unrestricted notebook runner. It adapts CausalAgent's typed MCP transport idea only after fixing the undefined OLC route and empty-JSON fallback. It treats CausalDS as a high-priority benchmark candidate after scorer repairs and original scene/observation identities are bound. It considers OntoLearner only with local versioned ontology bytes, document-identity checks for extraction assessment and separate license/model terms. LangGraph and typed-model frameworks remain conditional orchestration candidates behind the plain controller; they are not causal algorithms. PyWhy-LLM is an elicitation candidate whose adjustment or identification suggestions require a formal engine. Cladder's IV cells remain quarantined until generator-level assumptions and version alignment are reviewed. Unlicensed or source-paper/code-misaligned leads remain deferred.

The next experiment is an offline public or generated task set with a frozen brief, deterministic structured-interview baseline, fixed prompts and split, target/clock/diagram rubrics, malformed and refusal cases, resource budget and stopping rule. It should measure semantic preservation and adverse-case handling before any live backend call. A later live-model comparison needs a named provider, custody, price and independent adjudication; none is claimed here. Authentic bank truth, production reliability and customer economics remain later conditional evidence.


% SOURCE docs/release/2026-10-29/source/docs/writing-repair/2026-10-29/src/a0c37e73123737a5.tex CHAR 0
\section{An estimation assistant must preserve the target it validates}
\label{exec14:estimation-validation}

The choice of an estimator follows three earlier decisions: the business target,
the information available at the decision time, and identification under the
stated model. The assistant should keep these decisions separate from statistical
calibration and diagnostic adequacy. A balanced weighted sample can still omit
a common cause; an excellent predictive model can use a field recorded after
the intervention. A critic that approves the prose does not validate either
assumption. The causal co-pilot roadmap is valuable precisely because it makes
eligibility, baseline time, incident versus prevalent treatment, follow-up and
outcome definitions explicit before estimation \citep{v12_la_2407_19118}.

For an early arrears-review offer, let $X$ be the information available before
an eligible account's first qualifying arrears date, $A\in\{0,1\}$ the offer,
and $Y$ discounted net recovery over 180 days. The offer, acceptance and eventual
restructuring terms are distinct actions. The population $P$ consists of the
specified eligible accounts at time zero, with one declared interference and
sampling unit. Define $\tau_P=E_P\{Y(1)-Y(0)\}$. Offer cost is either included
in $Y$ or subtracted once when comparing actions. Changing the outcome horizon
or the eligible population creates a different target.

Under consistency, conditional exchangeability, positivity, a suitable sampling
law and the declared measurement conditions, write
$\mu_a(x)=E(Y\mid A=a,X=x)$ and $p(x)=\Pr(A=1\mid X=x)$. For candidate nuisance
functions $m_a$ and $e$, a population augmented inverse-probability score is
\[
 \phi=m_1(X)-m_0(X)+\frac{A\{Y-m_1(X)\}}{e(X)}
       -\frac{(1-A)\{Y-m_0(X)\}}{1-e(X)}.
\]
Conditioning on $X$ and subtracting the identified conditional contrast gives
\begin{equation}
 E(\phi\mid X)-(\mu_1-\mu_0)
 =(e-p)\left\{\frac{m_1-\mu_1}{e}
               +\frac{m_0-\mu_0}{1-e}\right\}.
 \label{exec14:aipw-remainder}
\end{equation}
This algebra explains the population double-robustness property when the
denominators are well defined: the conditional bias vanishes if the propensity
is right or both outcome regressions are right. It does not establish the causal
assumptions, a finite-sample confidence guarantee, or arbitrary learner validity.
Cross-fitting can separate nuisance training from score evaluation, but rates,
moments, dependence and uncertainty calculations still require their own
conditions. The inspected OCI notebook uses a cross-fitted AIPW approach within
a human-principal and actor/critic workflow; its source results remain distinct
from this project derivation \citep{v12_la_2607_22443,exec14_estimation_sources}.

\subsection{Full, trimmed and overlap targets answer different questions}

The exact three-stratum fixture in Table~\ref{exec14:target-table} gives a
concrete test of target preservation. Its probabilities, effects and monetary
interpretation are stipulated; no bank or package was fitted. The strata have
masses $p_s$, offer probabilities $e_s$ and conditional effects $\tau_s$.
\begin{table}[htbp]
\centering
\caption{Declared population law for a target-preservation fixture.}
\label{exec14:target-table}
\begin{tabular}{lrrr}
\toprule
Stratum & $p_s$ & $e_s$ & $\tau_s$\\
\midrule
Low offer propensity & $1/2$ & $1/20$ & $0$\\
Intermediate propensity & $3/10$ & $1/2$ & $4$\\
High offer propensity & $1/5$ & $19/20$ & $12$\\
\bottomrule
\end{tabular}
\end{table}
The full-population contrast is $\sum_s p_s\tau_s=18/5=3.6$. A fixed rule that
retains only propensities in $[.1,.9]$ keeps the intermediate stratum and yields
$4$. Conventional overlap weighting instead gives
\[
 \tau_{\mathrm{ov}}=
 \frac{\sum_s p_s e_s(1-e_s)\tau_s}{\sum_s p_s e_s(1-e_s)}
 =\frac{1656}{433}\simeq3.82448.
\]
These are three correct answers to three different questions. The target rule,
retained mass and weighting measure belong in the report next to the number.
Stabilizing an estimate by removing low-overlap accounts may be scientifically
reasonable, but it changes which accounts the recommendation concerns.

The pinned OCI notebook supplies an instructive source-level discrepancy. It
trims a propensity range, then applies \texttt{aipw\_summary} to the retained
pseudo-outcomes with uniform averaging while serializing that result as ATO.
The operation is a trimmed-population average; it is not automatically the
conventional $e(X)\{1-e(X)\}$-weighted overlap target. The proposed adapter must
rename the target with its rule or implement the intended weighting explicitly.
When trimming uses estimated propensities, the retained population is itself
estimated, adding selection and target uncertainty. The exact fixture above
uses known propensities and establishes only the arithmetic distinction
\citep{exec14_estimation_sources,exec14_estimation_oracles}.

The paper reports 231 ACIC datasets, a ten-dataset scaffolding comparison and
disguised LaLonde analyses. It reports better behavior with its structured
workflow and critic in those settings; these findings justify a candidate
architecture, not a claim of independently reproduced bank performance. A
critic's overlap refusal may be useful even when no effect is returned.
However, selecting critic-approved estimates does not automatically establish
coverage conditional on approval. Both the all-request and selected denominators
must be shown, with the selection rule fixed before final evaluation
\citep{v12_la_2607_22443}.

\subsection{Validation losses, unknown states and retrieval failure}

CausalDS separates narrative scenes, an observation layer and hidden causal
truth. Its identification targets refer to conceptual variables; an adapter must
not infer that a noisy observed proxy provides the same identifying law. The
paper's 100-scene exam, sampled from 953 scenes, reports different patterns for
symbolic success, numeric quality, abstention and resource cost. Those are dated
author findings for a benchmark. Its source scorer also exposes implementation
choices that matter when adapting the benchmark: missing identification metadata
can default toward an identifiable interpretation, and finite numeric endpoints
must be checked separately from successful conversion to a floating-point value
\citep{v12_la_2607_08093,exec14_estimation_sources}.

The interval score for a central interval $[l,u]$, target $\theta$ and nominal
error $\alpha$ is
\[
 S_\alpha(l,u;\theta)=(u-l)+\frac{2}{\alpha}(l-\theta)\mathbf1\{\theta<l\}
                  +\frac{2}{\alpha}(\theta-u)\mathbf1\{\theta>u\}.
\]
It penalizes width and noncoverage. The inspected CausalDS source reports both
uncapped and capped normalized scores, together with coverage and width. It is
incorrect to describe it as offering only capped loss. A dashboard that keeps
only the capped value can nevertheless reverse a comparison. Score pairs
$(1,100)$ and $(9,9)$ have uncapped means $50.5$ and $9$; with a cap of ten,
the means are $5.5$ and $9$. Such pairs can be constructed with covered
zero-centered intervals of the corresponding widths and truth zero, so the
arithmetic is compatible with the interval-score definition. This is an exact
illustration of a dashboard choice, not a result from executing CausalDS
\citep{exec14_estimation_sources,exec14_estimation_oracles}.

The proposed assessment keeps identified, nonidentified and unknown as separate
scientific states. Unsupported operations, malformed files, execution failures
and deliberate abstentions have separate operational states. It reports how
many requests returned finite valid intervals and how many were refused, while
retaining uncapped losses for the valid interval population. Assigning an
infinite or arbitrary loss to every refusal may be inappropriate for a task
where refusal is correct; silently dropping failures is also inappropriate.
Task-specific scientific scoring and the all-request operating denominator
must therefore be published together.

ARGUS provides a complementary example at the assumption-audit boundary. It
audits reported support for identification assumptions, not causal truth. Its
synthetic flaw results show that retrieval gating changes the tradeoff between
flagging flaws and missing omissions. In its five-paper human pilot, weak
retrieval produces excessive severity, and the reported adjustment is in-sample.
The useful remedy is to retain retrieval quality and evidence status as separate
fields, with unresolved evidence labeled unknown. A failure to retrieve a bank
policy is not evidence that the policy does not exist
\citep{v12_la_2609_30867}.

The initial prototype should adapt OCI's specification and critic artifacts
behind fixed bounded operations, and adapt CausalDS's hidden-truth and scoring
separation only after its generator and scoring paths are independently checked.
OCI's unbounded notebook-runner timeout is not an admissible default for the
bounded CLI/MCP service. Direct package integration, a live provider comparison
and bank validation are separate later claims. The concrete near-term output is
a reproducible design that preserves the target, refuses unsupported certainty,
and measures review burden against a deterministic analyst baseline.
\subsection{A workflow succeeds by preserving a decision}

The developed examples give each handoff a substantive test. A servicing rule
must react to its generated history; a rollout comparison must preserve cohort
composition and untreated-trend assumptions; an observation correction must
retain the population frame and measurement law. The architectural examples
then show how trimming, repeated criticism, failed tools and selective scoring
can change what a completed report means. These failures require different
remedies and cannot be summarized by one success rate.

The proposed controller should preserve the question, unresolved premises and
evidence through those operations, with deterministic checks wherever their
conditions can be made explicit. Its advantage over a structured analyst
workflow remains an empirical question about usable answers and complete
review cost. The next chapter compares the tools that could implement this
restricted role, including the source, rights and execution obligations that
a plausible architecture alone cannot settle.
\label{exec14:estimation-validation-end}

